\documentclass[11pt]{article}
\usepackage{amsmath,amssymb,amsthm,geometry}
\usepackage[hidelinks]{hyperref}
\geometry{margin=1in}
\newtheorem{proposition}{Proposition}
\newtheorem{remark}{Remark}
\newcommand{\R}{\mathbb R}
\newcommand{\dd}{\,\mathrm d}
\newcommand{\proj}{\operatorname{proj}}
\newcommand{\av}[1]{\left\langle #1\right\rangle}
\newcommand{\Cov}{\mathcal C}
\newcommand{\Cross}{\mathcal B}
\usepackage{mathrsfs}
\newtheorem{theorem}{Theorem}
\title{Navier--Stokes source propagation and the finite coupled bridge\\
Original third-growth scales, exact pulse map, and complete residual}
\author{}\date{}
\begin{document}\maketitle
This reader contains the complete preceding operator, pulse, covariance,
source-transport and radial-moment proofs. The new parts retain the
original third-growth seed and evolving parent, prove all amplitude
derivative costs, identify the exact source clock and tangent frame,
construct an invertible map to the actual source homogeneous pulse,
and calculate the entire compact curl candidate and its force increment.

The released source theorem is attributed to the supplied manuscript;
this lane has not independently certified its whole proof. Its positive
swirl path and the exact negative reflected path retain their respective
forces. The distinct finite coupled candidate has the explicit residual
proved here. The companion 102-page reader contains its preceding
original-stage proofs. No infinite modified coupled construction or
uniform terminal force estimate is asserted by the finite pulse map.
\tableofcontents\clearpage
\part{Complete preceding bridge and actual source propagation}
\section{Exact operator bridge to the released Navier--Stokes profiles}
\label{nsbridge:nsb:operator}

This bounded comparison uses the locally saved released Navier--Stokes source
(the source pages are 7--15 and 24--27 for the profile, stress and pulse
construction). Its stated physical viscosity is denoted by $\nu_{\rm NS}>0$.
The Boussinesq lane variables and physical viscosity remain $d$, $\mu$ and
$\nu=\sqrt{A_0}$; none is identified with $\nu_{\rm NS}$.

Use right-handed cylindrical coordinates $(r,\vartheta,z)$ and define
\[
 y_1=z,\qquad y_2=\frac{r^2}{2}>0,\qquad
 v=(u^z,ru^r)=J\nabla_y\Psi,
 \quad J=\begin{pmatrix}0&-1\\1&0\end{pmatrix}.
\]
The meridional divergence identity is
$\partial_{y_1}v_1+\partial_{y_2}v_2=0$. Define the circulation and
meridional vorticity variables
\[
 \Gamma=ru^\vartheta,
 \qquad \xi=\frac{\omega^\vartheta}{r},
 \qquad h=\frac{\Gamma}{2y_2}=\frac{u^\vartheta}{r},
 \qquad Q=\partial_{y_1}^2+2y_2\partial_{y_2}^2,
 \qquad M=Q+4\partial_{y_2}.
\]
The azimuthal momentum and meridional vorticity rows of the cylindrical
equations (with arbitrary smooth force components) become, on $y_2>0$,
\begin{align}
 D_v\Gamma&=\nu_{\rm NS}Q\Gamma+r f^\vartheta,
 \label{nsbridge:nsb:Gamma}\\
 D_v\xi&=\nu_{\rm NS}M\xi+\partial_{y_1}(h^2)
       +\operatorname{curl}_y(f^z,f^r/r),
 \label{nsbridge:nsb:xi}
\end{align}
where $D_v=\partial_t+v\cdot\nabla_y$ and
$\operatorname{curl}_y(a,b)=\partial_{y_1}b-\partial_{y_2}a$.

Indeed, $\partial_r=r\partial_{y_2}$ and
$\partial_r^2+r^{-1}\partial_r=2y_2\partial_{y_2}^2+2\partial_{y_2}$.
For $\Gamma$, the azimuthal viscous operator is
$\partial_r^2-r^{-1}\partial_r+\partial_z^2=Q$.
For $r\xi=\omega^\vartheta$, the cylindrical operator
$(\partial_r^2+r^{-1}\partial_r-r^{-2}+\partial_z^2)(r\xi)$,
divided by $r$, is $M\xi$; explicitly the radial terms are
$2y_2\xi_{y_2y_2}+4\xi_{y_2}$. The centrifugal term is
$r^{-2}\partial_z((u^\vartheta)^2)=\partial_{y_1}(\Gamma^2/(4y_2^2))
=\partial_{y_1}(h^2)$. These calculations prove every coefficient in
(\ref{nsbridge:nsb:Gamma})--(\ref{nsbridge:nsb:xi}), including the $r$ force conversion.

The velocity--vorticity relation is also retained. Since
$u^z=-\Psi_{y_2}$ and $u^r=\Psi_{y_1}/r$,
\[
 \xi=\frac{\partial_z u^r-\partial_r u^z}{r}
     =\frac{\Psi_{y_1y_1}}{2y_2}+\Psi_{y_2y_2}
     =\frac{Q\Psi}{2y_2}.
\]
Thus the two transport rows do not by themselves replace the elliptic
velocity reconstruction. Conversely, on a simply connected meridional
domain, a smooth $\Psi,\Gamma$ satisfying this relation and both transport
rows recovers $u$. Define its momentum residual without pressure as
$a=\partial_tu+(u\cdot\nabla)u-\nu_{\rm NS}\Delta u-f$.
The circulation row is $a^\vartheta=0$, while the vorticity row is
$\partial_z a^r-\partial_r a^z=0$. The one-form
$a^r\,\mathrm dr+a^z\,\mathrm dz$ is therefore closed, and
$p=-\int(a^r\,\mathrm dr+a^z\,\mathrm dz)$ along a path from a fixed
base point is path independent. Its derivatives are $p_r=-a^r$,
$p_z=-a^z$. This proves the local momentum reconstruction, with pressure
unique up to a function of time; boundary data are not removed.

The change from circulation to swirl rate is invertible on the exact retained
domain: $\Gamma=2y_2h$. Product differentiation gives
\[
 Q(2y_2h)=2y_2Mh,\qquad
 D_v(2y_2h)=2y_2D_vh+2v_2h.
\]
Therefore (\ref{nsbridge:nsb:Gamma}) is equivalent to
\begin{equation}
 D_vh=\nu_{\rm NS}Mh-\frac{v_2}{y_2}h+\frac{f^\vartheta}{r}.
 \label{nsbridge:nsb:h}
\end{equation}
The source profile uses precisely this relation: in its similarity variables,
$E$ is the azimuthal velocity profile and
$H_{\rm src}=\sqrt{2X}E$ is the scaled circulation profile (source equations
(4.3), (4.8)). Since $r=\sqrt{2qX}$ and
$u^\vartheta=q^{-A}E$, its physical circulation is
$\Gamma=q^{1/2-A}H_{\rm src}=q^{-h_{\rm src}}H_{\rm src}$,
where the source exponents satisfy $A=1/2+h_{\rm src}$. Here $q>0$ is the
geometric similarity coordinate and $h_{\rm src}>0$ is an exponent; neither
is a signed amplitude scale. This is the exact scale dictionary; it does not
identify $H_{\rm src}$ with the Boussinesq scalar amplitude.

The sign of the blowout amplitude must be kept separate from that positive
geometric coordinate. The axisymmetric equations have the exact involution
\[
 (u^r,u^\vartheta,u^z,p;f^r,f^\vartheta,f^z)
 \longmapsto
 (u^r,-u^\vartheta,u^z,p;f^r,-f^\vartheta,f^z).
 \label{nsbridge:nsb:sign}
\]
It sends $(\Gamma,h)$ to $(-\Gamma,-h)$, leaves $\xi$ and the meridional
flow unchanged, and preserves every displayed equation: the only quadratic
swirl term is $\partial_{y_1}(h^2)$, while the azimuthal equation is linear in
$(\Gamma,f^\vartheta)$. Consequently a source-selected positive profile
$E>0$ (source pp. 7--9) does not justify assuming that every blowout branch is
positive; its reflected branch has $E<0$, $\Gamma<0$ and the same blowup
magnitude. The exponent $q^{-h_{\rm src}}$ controls the magnitude in either
branch. The source's positive covariance coefficients and viscosity are
separate construction choices, not a sign claim about the velocity amplitude.

For comparison with the lane's controlled scalar equation, set
$B=h^2$. The globally smooth square equation, including its geometric and
viscous terms, is
\begin{equation}
 D_vB=\nu_{\rm NS}MB-2\nu_{\rm NS}
 |\nabla h|_Q^2-\frac{2v_2}{y_2}B+\frac{2h}{r}f^\vartheta,
 \qquad |\nabla h|_Q^2=h_{y_1}^2+2y_2h_{y_2}^2.
 \label{nsbridge:nsb:B}
\end{equation}
This follows by multiplying (\ref{nsbridge:nsb:h}) by $2h$ and using
$M(h^2)=2hMh+2|\nabla h|_Q^2$. At zeros of $h$ equation
(\ref{nsbridge:nsb:B}) remains smooth; dividing by $B$ is unnecessary and would lose
that domain. On a component where $B>0$, the equivalent expression is
$D_vB=\nu_{\rm NS}MB-\nu_{\rm NS}|\nabla B|_Q^2/(2B)
-2v_2B/y_2+2hf^\vartheta/r$.

The exact relation to the Boussinesq lane is consequently an operator
correspondence, not a scalar conjugacy. The lane's retained-viscosity scalar
row has $D_v\Theta=d\Delta\Theta+f_\Theta$ in Cartesian material
coordinates. A direct identification $\Theta=B$ would require, on the same
physical domain, both $M=\Delta$ after the nonlinear coordinate pullback and
$-2v_2B/y_2-2\nu_{\rm NS}|\nabla h|_Q^2+2hf^\vartheta/r$ to be absorbed into
one prescribed force. The first condition fails as an operator identity:
$M=\partial_{y_1}^2+2y_2\partial_{y_2}^2+4\partial_{y_2}$ has variable metric
and drift, whereas the lane's $\Delta$ is the Euclidean Cartesian Laplacian.
The geometric, gradient and force terms must therefore be retained or
calculated in the particular source state. Nonconstant profiles alone
would not prove that such terms never cancel.
The exact bridge is (\ref{nsbridge:nsb:Gamma})--(\ref{nsbridge:nsb:B}) with the
full residual terms retained. It neither claims that the lane's controls are
the source pulses nor changes the source's pressure reconstruction.

For the meridional velocity, the source's pulse construction supplies the
following further exact local dictionary (source pp. 11--15 and 74--77). If
$w$ is a divergence-free pulse and $\pi$ its pressure increment, then
\[
 R(u_B+w,p_B+\pi)=R(u_B,p_B)+L_{u_B}(w,\pi)
 +\nabla\!\cdot(w\otimes w),
\]
\[
 L_{u_B}(w,\pi)=\partial_tw+(u_B\cdot\nabla)w
 +(w\cdot\nabla)u_B-\nu_{\rm NS}\Delta w+\nabla\pi,
 \qquad \nabla\cdot w=0.
\]
For the moving local phase with wavevector $n$ and transverse amplitude
$t_m$, the source equation is
\[
 t_m'+\mathcal Kt_m+m^2d_N t_m+i k m n\,\pi_m=-f_m,
 \qquad n\cdot t_m=0,
 \qquad d_N=\varepsilon k^2|n|^2,
\]
with the source's moving shear matrix
$\mathcal K=\begin{pmatrix}0&-2F&0\\2F+RF_R&0&0\\G_R&0&0\end{pmatrix}$.
The transverse projection uses
\[
 A_\Phi=-\mathcal K+
 \frac{n(n^T\mathcal K-n'^T)}{|n|^2}.
\]
These are exact source pulse equations, including viscosity and pressure
projection. Their quadratic covariance is chosen to cancel the annular
stress; the source explicitly retains curl, cutoff, and correction interactions.
The Boussinesq lane's two-dimensional pair and its full profile residuals have
an exact algebraic form, but no identification with this three-dimensional
pulse equation follows without an additional embedding that supplies the
source's cylindrical geometry, moving plane, covariance cone and correction
cycle. That embedding is not asserted here.

The current bridge therefore proves a precise morphism of variables, operators,
forces and pulse residuals and records the exact obstruction to a direct scalar
identification. It establishes the next programme datum while preserving the
unproved status of any singular limit or third return.


\sloppy

\section{Scope and source equations}
This note records one local, exact comparison.  It does not identify the
coupled Boussinesq stages with the released Navier--Stokes construction and
makes no claim about a global Navier--Stokes solution.  The source is the
released PDF \texttt{openai\_navier\_stokes.pdf}, SHA-256
\texttt{8c8a94ad9ac824c8b605b9827cadf7beaca48bd10b380de3cfc872a2c37afa81}.
The source statements used here are Section 7, equations (7.2)--(7.8),
(7.12)--(7.22), and Section 9, equations (9.1)--(9.2), on PDF pages 74--75,
77--81, and 101--103. In the source notation, on one lifted pulse rectangle,
\begin{equation}
 \Phi=p\theta+p_zZ/\varepsilon+x_0R-v(pF+p_zG),\qquad
 n=\nabla_*\Phi,
\end{equation}
with
\begin{equation}
 n=(x_0-v(pF_R+p_zG_R),\ p/R,\ p_z-\varepsilon v(pF_Z+p_zG_Z)).
 \label{nsbridge:eq:nsource}
\end{equation}
For a nonzero integer harmonic $m$, the exact principal pulse equations are
\begin{equation}
 n\cdot t_m=0,\qquad
 t_m'+\mathsf Kt_m+m^2d_Nt_m+ikm n\pi_m=-f_m,
 \qquad d_N=\varepsilon k^2|n|^2,
 \label{nsbridge:eq:sourcepulse}
\end{equation}
where a prime is the pulse-coordinate derivative and
\begin{equation}
 \mathsf K=\begin{pmatrix}0&-2F&0\\2F+RF_R&0&0\\G_R&0&0\end{pmatrix},\qquad
 \mathsf A_\Phi=-\mathsf K+\frac{n(n^T\mathsf K-(n')^T)}{|n|^2}.
 \label{nsbridge:eq:sourceoperator}
\end{equation}
The source pressure coefficient is exactly
\begin{equation}
 \pi_m=\frac{i}{km|n|^2}\bigl(n\cdot\mathsf Kt_m-(n')\cdot t_m+n\cdot f_m\bigr).
 \label{nsbridge:eq:sourcepressure}
\end{equation}
For $m=1$, a real homogeneous pulse $t_1$ has the source energy identity
(7.22):
\begin{equation}
 \frac12(|t_1|^2)'=-g\cdot\bigl(t_{1,r}(t_{1,\theta},t_{1,z})\bigr)
 -\varepsilon k^2|n|^2|t_1|^2,
 \label{nsbridge:eq:sourceenergy}
\end{equation}
where $g=(RF_R,G_R)$ in the tangential coordinates.  Thus the source pulse
retains both shear transfer and viscous dissipation.

\section{The coupled physical pair}
For one activated layer of the coupled calculation, put
\begin{equation}
 a_c=-\frac{J\zeta\cdot G_{<q}}{Kr^2},\qquad b_c=K\zeta_1,\qquad
 \lambda_c^\Theta=\kappa K^2r^2,\qquad \lambda_c^\Omega=\varepsilon_c K^2r^2.
\end{equation}
Its exact retained-viscosity pair, including optional scalar and vorticity
controls, is
\begin{equation}
 \dot y=C_cy+c_c,\qquad y=\binom{\Theta}{\Omega},\qquad
 C_c=\begin{pmatrix}-\lambda_c^\Theta&a_c\\b_c&-\lambda_c^\Omega\end{pmatrix},\quad
 c_c=\binom{c_\theta}{c_\omega}.
 \label{nsbridge:eq:coupledpair}
\end{equation}
Here $G_{<q}$, $\zeta$, and $r$ evolve with the inherited background and phase
ODEs; in particular no fixed-background substitution is made.  Let $v$ be a
source pulse coordinate on an interval and let $q_v=\dd t/\dd v$ describe
its physical-time parametrization. The forward source parametrization has
$q_v>0$; the algebra below also retains a negative orientation without
changing the amplitude sign. Define $y'(v)=q_v(C_cy+c_c)$.
Here $K,r$ are the coupled frequency and phase length, not the source
carrier $k$ and cylindrical radius $R$. The original $A_0,\lambda_0,\mu$
and time scale $\nu=\sqrt{A_0}$ are unchanged. The source chart parameter
$\varepsilon$ is not identified with either physical diffusivity.

\section{Exact local residual map}
Let $B(v)\in\R^{3\times2}$ be a $C^1$ tangent frame with full column rank and
$n(v)^TB(v)=0$, so its range is exactly $n(v)^\perp$.  Let $B_\ell(v)$ be any
left inverse, $B_\ell B=I_2$; then $B B_\ell$ acts as the identity on
$n^\perp$.  Let $L(v)\in\R^{2\times2}$ be any $C^1$ coefficient map and embed
the coupled pair into a transverse source amplitude by
\begin{equation}
 t(v)=B(v)L(v)y(v).
 \label{nsbridge:eq:embedding}
\end{equation}
For an embedding with a non-tangent frame the transversality condition would be
\begin{equation}
 n(v)^TB(v)L(v)y(v)=0.
 \label{nsbridge:eq:transversecondition}
\end{equation}
For the tangent frame used here, the condition holds for every $L,y$.  Set
\begin{equation}
 H(v)=B_\ell(v)\bigl(\mathsf A_\Phi(v)B(v)-B'(v)\bigr).
 \label{nsbridge:eq:Hdef}
\end{equation}
Then the exact projected source residual of the embedded coupled pair is
\begin{align}
 \mathcal E_{L,c}:=B_\ell\proj_{n^\perp}
 \left[t'-\mathsf A_\Phi t+m^2d_Nt+f_m\right]
 &=\left[L'+q_vLC_c-HL+m^2d_NL\right]y+q_vLc_c \\
 &\quad+B_\ell\proj_{n^\perp}f_m.
 \label{nsbridge:eq:residualmap}
\end{align}
In particular, with $L=I$ and no added source, the defect is the explicit
matrix expression
\begin{equation}
 \mathcal E_{I,0}=\bigl(q_vC_c-H+m^2d_NI_2\bigr)y.
 \label{nsbridge:eq:identitydefect}
\end{equation}
This gives a local residual relation with every actual coefficient retained.

\begin{proof}
Differentiate \eqref{nsbridge:eq:embedding}:
$t'=B'L y+B L'y+B L y'$.  Since $y'=q_v(C_cy+c_c)$,
differentiate $n^TB=0$ to obtain $n^TB'=-(n')^TB$. Also
$n^T\mathsf A_\Phi B=-(n')^TB$, directly from
\eqref{nsbridge:eq:sourceoperator}. Thus $\mathsf A_\Phi B-B'$ has transverse
range and equals $BH$. Applying $B_\ell$ to the transverse terms and
$B_\ell\operatorname{proj}_{n^\perp}$ to the force proves
\eqref{nsbridge:eq:residualmap}. More explicitly,
\begin{equation}
 n^T(t'-\mathsf A_\Phi t+m^2d_Nt)=0,\qquad
 n^T(t'-\mathsf A_\Phi t+m^2d_Nt+f_m)=n^Tf_m.
\end{equation}
The unprojected force need not be transverse. The pressure below removes
precisely this remaining normal component.
\end{proof}

\section{Pressure reconstruction and full three-dimensional residual}
For any embedded transverse $t$ and arbitrary $f_m$, define
\begin{equation}
 \pi_{L,c}=\frac{i}{km|n|^2}
 \left(n\cdot\mathsf Kt-(n')\cdot t+n\cdot f_m\right).
 \label{nsbridge:eq:pressuremap}
\end{equation}
Then the full source residual satisfies the exact decomposition
\begin{align}
 \mathcal R_{L,c}
 &: =t'+\mathsf Kt+m^2d_Nt+ikm n\pi_{L,c}+f_m \\
 &=\proj_{n^\perp}\left(t'+\mathsf Kt+m^2d_Nt+f_m\right)
 =B\,\mathcal E_{L,c},
 \label{nsbridge:eq:fullresidual}
\end{align}
where the last equality uses the definition of $H$ and the transverse condition.
Thus the pressure removes precisely the normal component; it does not remove
$\mathcal E_{L,c}$, the moving-frame, shear, diffusion, or coupled-background
mismatch.  If one elects to force the source pulse to obey the released equation,
then the unique additional transverse forcing is
\begin{equation}
 f_{m,\mathrm{corr}}=-B\,\mathcal E_{L,c},
\end{equation}
added to the old $f_m$. The old $\mathcal E_{L,c}$ is evaluated before
this addition, so the new projected residual is zero.
Since $n^Tf_{m,\mathrm{corr}}=0$, pressure is unchanged.
This is an exact local
forced relation, not an identification of the physical constructions.

\section{Diffusion and energy comparison}
The source diffusion multiplier is $m^2d_N=m^2\varepsilon k^2|n|^2$ in
\eqref{nsbridge:eq:sourcepulse}.  The coupled pair has the two retained physical losses
$\lambda_c^\Theta$ and $\lambda_c^\Omega$, which are generally unequal to
$m^2d_N$.  Their signed discrepancy in \eqref{nsbridge:eq:identitydefect} is explicitly
\begin{equation}
 \bigl(m^2d_N-q_v\lambda_c^\Theta\bigr)\Theta,
 \qquad
 \bigl(m^2d_N-q_v\lambda_c^\Omega\bigr)\Omega
\end{equation}
when $H=0$ and $C_c$ is expanded entrywise; with general $H$ the exact
diagonal entries are those of $q_vC_c-H+m^2d_NI_2$.  No damping is dropped.
If one compares nonnegative loss magnitudes instead, the corresponding cost
gaps are $|m^2d_N-q_v\lambda_c^\Theta|$ and
$|m^2d_N-q_v\lambda_c^\Omega|$.
The required finite-interval coordinate map can be constructed explicitly.
Let $I=[v_0,v_1]$, and let $X_c,X_s$ be the identity-initial-value
fundamental matrices for $D_c=q_vC_c$ and $D_s=H-m^2d_NI_2$,
respectively. For either continuous matrix $D$, define
\[
 X_D(v)=I_2+\sum_{j\ge1}
 \int_{v_0\le w_j\le\cdots\le w_1\le v}
 D(w_1)\cdots D(w_j)\,\dd w_j\cdots\dd w_1 .
\]
With $M_D=\sup_I\|D\|$ and $T_I=v_1-v_0$, the $j$th integral has norm
at most $(M_DT_I)^j/j!$. Uniform convergence proves the integral equation
$X_D=I_2+\int_{v_0}^vDX_D$ and hence $X_D'=DX_D$.
The same iteration for $Z'=-ZD$, $Z(v_0)=I_2$ gives $ZX_D=I_2$
by differentiation, proving invertibility on $I$. For a specified
invertible $L_0$, set
\[
 L=X_sL_0X_c^{-1},\qquad L^{-1}=X_cL_0^{-1}X_s^{-1}.
\]
This has norm at most $\|L_0\|e^{(M_c+M_s)T_I}$; its inverse has the
same bound with $L_0^{-1}$. Differentiation proves
\begin{equation}
 L'=HL-q_vLC_c-m^2d_NL,
 \label{nsbridge:eq:intertwiner}
\end{equation}
so the homogeneous embedded pair has zero projected residual. This proves
an actual finite-interval conjugacy along the retained coefficient path.
It does not make the source phase,
pressure, cutoff, or nonlinear covariance equal to the coupled Boussinesq fields.
For general controls $c_c$, the forced term is $q_vLc_c$ in
\eqref{nsbridge:eq:residualmap}.

For an arbitrary embedded field (possibly complex), the exact energy balance
retains its full defect:
\begin{equation}
 \begin{split}
 \frac12(|t|^2)'={}&-\operatorname{Re}(\overline t^{\,T}\mathsf Kt)
 -m^2d_N|t|^2-\operatorname{Re}(\overline t^{\,T}f_m)\\
 &+\operatorname{Re}(\overline t^{\,T}B\mathcal E_{L,c}).
 \end{split}
\end{equation}
This follows by taking the real Hermitian product of
\eqref{nsbridge:eq:fullresidual} with $t$. Its pressure term vanishes because $n$ is
real and $n^Tt=0$. The $2F$ part of $\mathsf K$ is skew-symmetric, hence
its real work is zero; the shear work equals
$\operatorname{Re}(\overline{t_r}(g_\theta t_\theta+g_z t_z))$.
Only for a zero residual, zero force, $m=1$, and real $t$ does this reduce
to \eqref{nsbridge:eq:sourceenergy}. The real coupled pair has
\begin{equation}
 \frac12\frac{\dd}{\dd t}|y|^2
 =-\lambda_c^\Theta\Theta^2-\lambda_c^\Omega\Omega^2
 +(a_c+b_c)\Theta\Omega+c_\theta\Theta+c_\omega\Omega,
 \label{nsbridge:eq:coupledenergy}
\end{equation}
which is a two-component physical energy identity.  Equations
\eqref{nsbridge:eq:residualmap} and \eqref{nsbridge:eq:coupledenergy} exhibit exactly where the
source's three-dimensional shear flux and the coupled pair's two-dimensional
exchange differ.

\section{Proven scope}
The source uses a cylindrical, localized, divergence-free curl construction
(Section 7.4 and Section 9, PDF pages 85--87 and 102--103); the displayed
principal pulse equation is only one component of that construction.  The
comparison above proves an exact local map and its full pressure-projected
residual for every finite interval on which the denominators and frame left
inverse exist.  It preserves source viscosity, moving phase normal, pressure,
rotation/shear matrix, and all coupled physical diffusion coefficients.  It does
not prove that the coupled third-growth or steering stages produce the source's
pulse phase $\Phi$, covariance, cutoff tails, or smooth terminal force.  Those
remain separate bridge problems.

The sign conventions in this comparison are also explicit.  The quantities
$q_v>0$, $\varepsilon>0$, and $d_N=\varepsilon k^2|n|^2>0$ record the chosen
orientation of physical time and the positive viscous coefficient; they are
not assumptions on the sign of a pulse or blowout amplitude.  The equations
\eqref{nsbridge:eq:sourcepulse}--\eqref{nsbridge:eq:sourcepressure} are linear in $(t_m,f_m)$,
so the simultaneous replacement $(t_m,f_m,\pi_m)\mapsto
(-t_m,-f_m,-\pi_m)$ is an exact reflected branch.  Likewise, the local map
\eqref{nsbridge:eq:embedding} is sign-equivariant under $y\mapsto-y$ and
$c_c\mapsto-c_c$, with the coefficients fixed at the actual path.
This is a statement about the retained linear equations, not a sign symmetry
of the entire coupled nonlinear system.
The source covariance is quadratic, so reversing a wave leaves that
covariance unchanged. Reversing the full Navier--Stokes velocity need not
reverse its force, since $(u\cdot\nabla)u$ is even in $u$.
The full-flow map that reverses axisymmetric swirl is a spatial reflection;
its complete formula is proved in the accompanying covariance bridge.
No sign of a blowout or endpoint is inferred from positive geometric scales.


\section{Source objects and the finite calculation}
This calculation uses the preserved 165-page supplied manuscript
\emph{Finite Time Blowup for Navier--Stokes}, source equations
(7.23)--(7.42), (9.12)--(9.14). Its PDF SHA-256 is
\begin{center}\small
\texttt{8c8a94ad9ac824c8b605b9827cadf7beaca48bd10b380de3cfc872a2c37afa81}.
\end{center}
The source pulse fields, auxiliary rectangles, phases, and leading
covariance columns are the given objects; their global construction is
source input. Every finite map asserted below is proved here. The existing
coupled-stage proofs remain in the complete 102-page reader. The companion
pulse note proves the exact frame residual and a constructed finite-interval
intertwiner along its evolving coefficients. The source uses positive
$q(z,t)$, band scales $Q_\beta=2^{-\ell_\beta}$,
$\varepsilon_\beta=Q_\beta^h$, $A=1/2+h$; none is set to one here or
identified with the original coupled $A_0,\lambda_0,\mu,\nu=\sqrt{A_0}$
or its physical diffusivities.

All operations in this note are on a compact set inside $q>0$ and the
source active annulus. There are finitely many relevant bands and
locally finite source labels. Supports of newly prescribed inputs are
strictly inside the annulus, so no lower bound on a leading amplitude
at a vanishing shell edge is assumed. Endpoint and shell-edge estimates
are not inferred from compact-interior estimates.

\section{Physical curl of the actual retained amplitude}
Let $y=(\Theta,\Omega)^T$ be the actual smooth retained coupled solution,
$y'=q_v(C_cy+c_c)$, with the coefficients of the companion pulse note.
Use the smooth actual coefficients on the chosen slow/auxiliary patch;
all harmonic coefficients and cutoffs are independent of $\theta$ in
cylindrical components, and each carrier has nonzero integer angular
frequency. For its source tangent frame $B$, smooth coefficient map $L$,
and a fixed smooth real compact cutoff $\chi$ of this type, put
\[
 t_m=\chi BLy,\qquad
 \mathcal E_\chi=
 \chi\{[L'+q_vLC_c-HL+m^2d_NL]y+q_vLc_c\}+\chi'Ly,
 \quad H=B_\ell(\mathsf A_\Phi B-B').
\]
Here prime differentiates along the pulse coordinate; the other slow
and auxiliary derivatives remain in the spatial formulas below.
Smooth parameter dependence of the companion's constructed $L$ follows
directly from its integral equations on a fixed interval: for a nonzero
parameter multi-index $I$,
\[
 \partial^I X(v)=
 \int_{v_0}^v\left[D\,\partial^I X+
 \sum_{0<J\le I}\binom{I}{J}(\partial^JD)
                                 \partial^{I-J}X\right]\,\dd w .
\]
The initial parameter derivatives vanish for identity initial data.
Inductively the known forcing is smooth and bounded on the compact patch;
the same convergent ordered-integral iteration solves this equation.
This proves every finite mixed derivative of $X,X^{-1},L$.
Since $n^TB=0$, $n^Tt_m=0$. Differentiation proves
$t_m'-\mathsf A_\Phi t_m+m^2d_Nt_m=B\mathcal E_\chi$.
Thus even the pulse-cutoff term is explicit.

A physical version of the exact curl construction makes the intermediate
map to momentum residual unambiguous. For each chosen harmonic let
$\phi(x,t)=k_\beta m\Phi_\beta(x,t,Y(x,t))$ be its real evaluated phase,
let $\xi=\nabla_x\phi\ne0$, and let $a(x,t)=Q_\beta^{-A}t_m$ be the
physical transverse amplitude. The source coordinate change is
$r=\sqrt{Q_\beta}R$, $z=Q_\beta^D Z$, $1-t=Q_\beta T$,
with $D=1/2-h$ and normalized axial derivative
$D_z=\varepsilon_\beta\partial_Z$. It gives
$\xi=(k_\beta m/\sqrt{Q_\beta})n_\Phi$; all auxiliary chain derivatives
are included in $n_\Phi$. Hence $\xi\cdot a=0$. Define in Cartesian
components
\begin{equation}
 C=\frac{i\,\xi\times a}{|\xi|^2},\quad
 \mathcal A=C e^{i\phi},\quad
 \nabla\times\mathcal A=(a+r_a)e^{i\phi},\qquad
 r_a=\nabla\times C. \label{nsbridge:cb:curl}
\end{equation}
The vector triple-product identity gives
$i\xi\times C=a$, proving the last equality. It also shows that the
physical potential has exactly the source factor $Q_\beta^{1/2-A}$
relative to its chart coefficient
$i\,n_\Phi\times t_m/(k_\beta m|n_\Phi|^2)$.
Cartesian differentiation of $C$ includes derivatives of the cylindrical
basis; the same terms are the $R^{-1}$ terms of source (7.37)--(7.38).
Take real parts and sum the finite retained harmonics. This constructs
a smooth compact velocity $w_y$ with $\nabla\cdot w_y=0$.

Retain a finite smooth source state $(u_*,p_*)$, including its background,
actual leading waves and their curl corrections, and the real
physical pressure increments obtained from the companion pulse pressure
formula with $f_m=0$:
\[
 \pi_{\chi,m}=\frac{i}{k_\beta m|n|^2}
                  \bigl(n^T\mathsf Kt_m-(n')^Tt_m\bigr),\qquad
 \pi_y=\sum_{\beta,m}\operatorname{Re}
             (Q_\beta^{-2A}\pi_{\chi,m}e^{i\phi_{\beta,m}}).
\]
These are smooth and compact wherever $t_m$ is. This definition fixes the
pressure map without assuming the pulse residual vanishes.
Define the actual candidate residual
\begin{align}
 R_y={}&R_{\nu_{\rm NS}}(u_*+w_y,p_*+\pi_y), \nonumber\\
 R_{\nu_{\rm NS}}(u,p)
   :={}&\partial_tu+(u\cdot\nabla)u-\nu_{\rm NS}\Delta u+\nabla p .
 \label{nsbridge:cb:Ry}
\end{align}
This is a definition by explicit differentiated fields, not an assumed
source term. Its exact expansion is
\[
 R_y=R_{\nu_{\rm NS}}(u_*,p_*)
 +\partial_tw_y+(u_*\cdot\nabla)w_y+(w_y\cdot\nabla)u_*
 -\nu_{\rm NS}\Delta w_y+\nabla\pi_y
 +(w_y\cdot\nabla)w_y .
\]
Every cutoff, parent, phase, and curl derivative occurs in this formula.
Physical $\nu_{\rm NS}$ is retained and is distinct from every source
time-scaling convention. Define
$\Delta R_y=R_y-R_{\nu_{\rm NS}}(u_*,p_*)$.
We correct the new compact increment; this difference has compact support
since every added term does. The baseline residual remains explicitly.

The averaging needed next is the source extended-field operation, not a
spatial average of the evaluated field. Evaluate derivatives on the
extended coefficients first by the complete chain rule. In a fixed
chart set $\mathcal D_i=\partial_{x_i}
+\sum_\alpha(\partial_{x_i}Y_\alpha)\partial_{Y_\alpha}$ and similarly
$\mathcal D_t$. These are the pullbacks of commuting physical
derivatives, so $[\mathcal D_i,\mathcal D_j]=0$ on their lifted domain.
For smooth periodic coefficients their auxiliary derivative terms have
zero Haar average. In cylindrical components angular differentiation
also differentiates the basis. This defines an extended residual
$\widetilde R_y$ whose evaluation is exactly \eqref{nsbridge:cb:Ry}.
Average the tangential components of the increment at fixed slow variables:
\begin{equation}
 F_2=\av{\Delta\widetilde R_{y,\theta}},\qquad
 F_1=\av{\Delta\widetilde R_{y,z}},\qquad
 \av{g}=\int_{\mathbb T^2}\frac1{2\pi}
                    \int_0^{2\pi}g\,\dd\theta\,\dd Y . \label{nsbridge:cb:F}
\end{equation}
These $F_e$ are smooth and compact
in the chosen physical annulus. Equations \eqref{nsbridge:cb:curl}--\eqref{nsbridge:cb:F}
are the explicit intermediate map from the retained coupled pair to a
signed momentum residual. They do not identify it with the original
two-dimensional Boussinesq force.

\section{Radial moments, support, and the original chart scales}
For each fixed $(z,t)$, choose $0<r_-<r_+$ containing the compact
supports of both $F_e$ and the chosen bumps in the active annulus.
The source bump is
\[
 b_e(r,z,t)=q^{-(e+1)/2}\widehat b_e(r/\sqrt q),\qquad
 \int_0^\infty x^e\widehat b_e(x)\,\dd x=1,\qquad e=1,2,
\]
with interior support. The change $r=\sqrt q\,x$ proves
$\int r^eb_e\,\dd r=1$ exactly. Define
\begin{equation}
 M_e=\int_0^\infty r^eF_e(r)\,\dd r,\qquad
 \sigma_e(r)=-r^{-e}\int_0^r s^e(F_e(s)-b_e(s)M_e)\,\dd s .
 \label{nsbridge:cb:primitive}
\end{equation}
Then
\begin{equation}
 (\partial_r+e/r)\sigma_e=-F_e+b_eM_e,\qquad
 \int_0^\infty r^e(F_e-b_eM_e)\,\dd r=0. \label{nsbridge:cb:radial}
\end{equation}
Differentiate $r^e\sigma_e$ to prove the first identity; integrate and
use the bump normalization for the second. The integral is zero below
$r_-$ and above $r_+$, so $\sigma_e$ has compact smooth support in the
same enclosing interval. A homogeneous solution is $cr^{-e}$, and a
compact homogeneous solution has $c=0$; thus the primitive is unique
among compact solutions of its displayed equation. Moreover a compact
solution of $(\partial_r+e/r)\sigma=-F_e$ exists if and only if $M_e=0$:
necessity follows by integrating the weighted derivative, and
sufficiency follows from \eqref{nsbridge:cb:primitive}. The two moments are
therefore the exact surviving finite-dimensional terms.

For $C_e=\int_{r_-}^{r_+}s^e\,\dd s$,
$|M_e|\le C_e\|F_e\|_\infty$ and
\[
 \|\sigma_e\|_\infty\le r_-^{-e}C_e
        (\|F_e\|_\infty+\|b_e\|_\infty|M_e|).
\]
If $G_e=F_e-b_eM_e$, successive radial derivatives obey the exact recurrence
\[
 \partial_r^{j+1}\sigma_e=-\partial_r^jG_e
 -e\sum_{i=0}^j \binom{j}{i}(-1)^ii!r^{-i-1}
                              \partial_r^{j-i}\sigma_e .
\]
This proves every finite radial derivative bound inductively. Parameter
derivatives follow by differentiating the compact integral; all derivatives
of $q,b_e$ and the moving physical data are retained. On the present compact
set $q$ is bounded away from zero. No terminal uniform estimate is implied.

The source (9.12)--(9.14) uses physical
$F_e=Q^{-2A-1/2}E_e^*$, where
$E_2^*=\av{E_\theta}_Y$ and $E_1^*=\av{E_z}_Y$.
For our fields define $E_e^*=Q^{2A+1/2}F_e$ in that same convention.
With $r=\sqrt Q R$ and $Q$ fixed in a chart, the exact conversions are
\begin{gather}
 M_e=Q^{e/2-2A}M_e^*,\quad M_e^*=\int R^eE_e^*\,\dd R,\quad
 b_e^*(R)=Q^{(e+1)/2}b_e(\sqrt Q R),\nonumber\\
 \Sigma_e(R)=Q^{2A}\sigma_e(\sqrt Q R),\qquad
 (\partial_R+e/R)\Sigma_e=-E_e^*+b_e^*M_e^*. \label{nsbridge:cb:chart}
\end{gather}
Substitute $r=\sqrt Q R$ including its measure in the moment, and
differentiate $\Sigma_e$ to prove each factor. The input $F_e$ is already
auxiliary independent. Taking only an averaged moment of an unaveraged
$F(R,Y)$ would not suffice: $F=b_e(R)\cos(2\pi Y_1)$ has zero averaged
moment but a nonzero exterior primitive at each generic $Y$.

\section{The source covariance derivative and its signed right inverse}
At each slow point let $w_{\rm tan}=(w_\theta,w_z)$ and define
\[
 \Cov(w)=\av{w_rw_{\rm tan}},\qquad
 \Cross(w,v)=\av{w_rv_{\rm tan}+v_rw_{\rm tan}}.
\]
Direct expansion proves
$\Cov(w+v)=\Cov(w)+\Cross(w,v)+\Cov(v)$ and
$D\Cov(w)[v]=\Cross(w,v)$.

Use the exact source real pulses
$b_\sigma=\chi_g(\xi_g)\psi(v)t_\sigma^h\cos(k\Phi_\sigma)$,
$\sigma\in\{+,-\}$. They exclude the slow cutoff $\eta_\beta$ and
scalar amplitudes; $\chi_g$ is an auxiliary transverse cutoff.
The two signs have disjoint auxiliary supports. Put
$H_{\rm cov}=[\Cov(b_+)\mid\Cov(b_-)]$.
For slow scalar coefficients independent of both averages,
\[
 \Cov(a_+b_++a_-b_-)=H_{\rm cov}(a_+^2,a_-^2)^T.
\]
All cross products vanish by disjointness. The factor $1/2$ from
$\av{\cos^2(k\Phi_\sigma)}_\theta$ is already inside the columns.

On the source interior cone the fixed primary amplitudes satisfy
$a_\sigma=\sqrt{(H_{\rm cov}^{-1}T_{0,*})_\sigma}>0$,
$T_{0,*}=(Q/q)^{A+1/2}T_0$, and
$W_0=\sqrt\varepsilon\sum_\sigma a_\sigma b_\sigma$.
For any signed slow real vector $\Sigma$, define
\begin{equation}
 d_\Sigma=H_{\rm cov}^{-1}\Sigma/\varepsilon,\qquad
 \delta a_\sigma=(d_\Sigma)_\sigma/(2a_\sigma),\qquad
 L_\Sigma=\sqrt\varepsilon\sum_\sigma\delta a_\sigma b_\sigma .
 \label{nsbridge:cb:L}
\end{equation}
It follows exactly that
\begin{gather}
 \Cross(W_0,L_\Sigma)=\varepsilon H_{\rm cov}(2a\odot\delta a)=\Sigma,
 \label{nsbridge:cb:inverse}\\
 \Cov(W_0+L_\Sigma)
   =\varepsilon T_{0,*}+\Sigma+
       \varepsilon H_{\rm cov}(\delta a_+^2,\delta a_-^2)^T.
 \label{nsbridge:cb:quad}
\end{gather}
This proves a linear right inverse of the derivative. It is not an
exact inverse of the quadratic covariance map. There is no sign or
size restriction on $\Sigma$ in these identities; no square root of
$H_{\rm cov}^{-1}(T_{0,*}+\Sigma/\varepsilon)$ is taken.

With Euclidean and induced operator norms and
$a_{\min}=\min_\sigma|a_\sigma|>0$, they give the explicit bound
\begin{equation}
 \|\Cov(L_\Sigma)\|\le
 \frac{\|H_{\rm cov}\|\|H_{\rm cov}^{-1}\|^2}
              {4\varepsilon a_{\min}^2}\,\|\Sigma\|^2. \label{nsbridge:cb:bound}
\end{equation}
Indeed $\|\delta a\|\le\|H_{\rm cov}^{-1}\|\|\Sigma\|/
(2\varepsilon a_{\min})$ and
$\|(\delta a_\sigma^2)_\sigma\|_2\le\|\delta a\|_2^2$.
All higher derivative costs are explicit from Leibniz and
\[
 \partial^I H_{\rm cov}^{-1}
 =-H_{\rm cov}^{-1}
 \sum_{0<J\le I}\binom{I}{J}(\partial^JH_{\rm cov})
                           \partial^{I-J}H_{\rm cov}^{-1},
\]
and, for any nonzero scalar $a$,
$\partial^I(a^{-1})=-a^{-1}\sum_{0<J\le I}\binom{I}{J}
(\partial^Ja)\partial^{I-J}(a^{-1})$.
Together with \eqref{nsbridge:cb:L} these recurrences give finite bounds for
each order without assuming constant amplitudes or constant columns.

\section{Assembly, the actual covariance change, and residual cancellation}
For the physical target $\sigma=(\sigma_2,\sigma_1)$ from
\eqref{nsbridge:cb:primitive}, use the source squared slow partition
$\sum_\beta\eta_\beta^2=1$ on its support. Set
\begin{align}
 W_{0,\rm phys}&=\sum_\beta Q_\beta^{-A}\eta_\beta W_{0,\beta},\nonumber\\
 v_{\rm pr}&=\sum_\beta Q_\beta^{-A}\eta_\beta
                        L_\beta(Q_\beta^{2A}\sigma).
 \label{nsbridge:cb:assembly}
\end{align}
Overlapping slow labels have disjoint auxiliary supports. Taking cross
covariance and using \eqref{nsbridge:cb:inverse} gives
\[
 \Cross(W_{0,\rm phys},v_{\rm pr})
 =\sum_\beta Q_\beta^{-2A}\eta_\beta^2Q_\beta^{2A}\sigma=\sigma.
\]
Every band scale remains distinct. Likewise the leading physical factor is
$Q^{-2A}\varepsilon(Q/q)^{A+1/2}
=Q^{-A+h+1/2}q^{-A-1/2}=q^{-A-1/2}$ because $A=1/2+h$.

Apply the exact curl \eqref{nsbridge:cb:curl} to each coefficient of
$v_{\rm pr}$, including the slow cutoffs before curling. Write the
resulting divergence-free increment as $v_s=v_{\rm pr}+r_s$.
Let $t_{s,m}$ denote the chart transverse harmonic coefficients of
$\eta_\beta L_\beta(Q_\beta^{2A}\sigma)$, before multiplication by
$Q_\beta^{-A}$. For each such coefficient choose
$\pi_{s,m}=i(n^T\mathsf Kt_{s,m}-(n')^Tt_{s,m})/
(k_\beta m|n|^2)$ and sum
$\operatorname{Re}(Q_\beta^{-2A}\pi_{s,m}e^{i\phi_{\beta,m}})$
to define the smooth compact $\pi_s$. Its full gradients remain below.
Let the actual old wave $w$ consist of the wave part of $u_*$
plus the candidate $w_y$, including all existing curl corrections, and write
$w=W_{0,\rm phys}+e_w$. Define tensor covariance
$\mathcal T(w)=\av{w\otimes w}$ and
$\mathcal B_T(u,v)=\av{u\otimes v+v\otimes u}$.
Expansion, with $e_w=w-W_{0,\rm phys}$, proves
\begin{equation}
 \Delta\mathcal T=
 \mathcal B_T(W_{0,\rm phys},v_{\rm pr})
 +\mathcal B_T(e_w,v_{\rm pr})
 +\mathcal B_T(w,r_s)+\mathcal T(v_s).
 \label{nsbridge:cb:fullcov}
\end{equation}
Its first radial--tangential pair is $\sigma$.
All other entries and the remaining three tensors are retained.
For example their radial--tangential pair has bound
\[
 2\|e_w\|_{L^2_{\theta,Y}}\|v_{\rm pr}\|_{L^2_{\theta,Y}}
 +2\|w\|_{L^2_{\theta,Y}}\|r_s\|_{L^2_{\theta,Y}}
 +\|v_s\|_{L^2_{\theta,Y}}^2 ,
\]
by Cauchy--Schwarz. The compact-set estimates above and the explicit
curl differentiate to bound all these terms at each finite order.

For clarity the complete residual after adding $v_s$ and its selected
pressure $\pi_s$ is, without any averaging,
\begin{align}
 R_{\nu_{\rm NS}}(u+v_s,p+\pi_s)
 =R_{\nu_{\rm NS}}(u,p)
 &+\partial_tv_s+(u\cdot\nabla)v_s+(v_s\cdot\nabla)u \nonumber\\
 &-\nu_{\rm NS}\Delta v_s+\nabla\pi_s
 +(v_s\cdot\nabla)v_s . \label{nsbridge:cb:residual}
\end{align}
This follows by expanding the product and retains viscosity and every
pressure term. For a symmetric tensor $S$ the averaged tangential
divergence in physical cylindrical coordinates is
\[
 (\av{\nabla\cdot S})_\theta
   =(\partial_r+2/r)\av{S_{r\theta}}+\partial_z\av{S_{z\theta}},
 \qquad
 (\av{\nabla\cdot S})_z
   =(\partial_r+1/r)\av{S_{rz}}+\partial_z\av{S_{zz}} .
\]
To verify, write the Cartesian divergence in the orthonormal cylindrical
basis, using $\partial_\theta e_r=e_\theta$,
$\partial_\theta e_\theta=-e_r$; periodic angular derivatives average
to zero and the two basis terms give $2/r$ and $1/r$ respectively.
The auxiliary chain terms average to zero by periodicity, but the
physical slow derivatives remain.

Let $U=u_*+w_y$, $P_y=p_*+\pi_y$, and set
\[
 \mathcal L=
 \partial_tv_s+(U\cdot\nabla)v_s+(v_s\cdot\nabla)U
 -\nu_{\rm NS}\Delta v_s+\nabla\pi_s .
\]
For a fully explicit averaged remainder, write $U=u_{\rm mean}+w$,
where $u_{\rm mean}$ is its angular mean, and put
$\mathcal V=\mathcal L-(w\cdot\nabla)v_s-(v_s\cdot\nabla)w$.
Then the average of \eqref{nsbridge:cb:residual} is
$\av{R_y}+\av{\mathcal V}+\nabla\cdot\Delta\mathcal T$.
Let $S_0=\mathcal B_T(W_{0,\rm phys},v_{\rm pr})$ and
$S_{\rm rem}=\mathcal B_T(e_w,v_{\rm pr})
 +\mathcal B_T(w,r_s)+\mathcal T(v_s)$.
Consequently the radial divergence of the first pair in
\eqref{nsbridge:cb:fullcov} changes the prescribed averaged residual $F_e$ into
$b_eM_e$ by \eqref{nsbridge:cb:radial}. The axial tensor divergences, the other
three tensors, the exact linear terms and the pressure terms of
\eqref{nsbridge:cb:residual} remain in the new residual. Equivalently the new
averaged tangential residual is exactly
\begin{align*}
 \av{R_{\rm new}}_{\rm tan}
 ={}&\av{R_{\nu_{\rm NS}}(u_*,p_*)}_{\rm tan}
 +(b_2M_2,b_1M_1)+\av{\mathcal V}_{\rm tan}\\
 &+\partial_z((S_0)_{z\theta},(S_0)_{zz})
   +(\nabla\cdot S_{\rm rem})_{\rm tan}.
\end{align*}
Every term here is defined by the displayed differentiated fields and
tensors. None is assumed to vanish.
The map from $y$ through \eqref{nsbridge:cb:F}, \eqref{nsbridge:cb:primitive},
\eqref{nsbridge:cb:L}, and \eqref{nsbridge:cb:assembly} is now an actual composite
construction. It retains the signed residual produced by the coupled
candidate and the source's original physical and chart units.

\section{Signs and the exact full-flow reflection}
The maps $F_e\mapsto M_e\mapsto\sigma_e\mapsto L_\Sigma$ are linear
and reverse sign with their inputs; $\Cov(L_\Sigma)$ is even.
The two primary amplitudes may also be replaced independently by
$s_\sigma a_\sigma$, $s_\sigma\in\{-1,1\}$, provided the same fixed
choice is used in each denominator: the cross identity is unchanged.
This does not negate the nonlinear Navier--Stokes equation.

Here is the exact full-flow map. In Cartesian coordinates let
$O=\operatorname{diag}(1,-1,1)$ and define
\[
 u^\sharp(x,t)=O u(Ox,t),\quad
 p^\sharp(x,t)=p(Ox,t),\quad f^\sharp(x,t)=O f(Ox,t).
\]
The chain rule gives
$\nabla u^\sharp=O(\nabla u)(Ox,t)O$,
$\Delta u^\sharp=O(\Delta u)(Ox,t)$, and
$(u^\sharp\cdot\nabla)u^\sharp
 =O((u\cdot\nabla)u)(Ox,t)$.
Thus $\nabla\cdot u^\sharp=0$ and
$R_{\nu_{\rm NS}}(u^\sharp,p^\sharp)=O R_{\nu_{\rm NS}}(u,p)(Ox,t)$
with the same positive viscosity.
In cylindrical components this is
\[
 (u_r^\sharp,u_\theta^\sharp,u_z^\sharp)(r,\theta,z,t)
 =(u_r,-u_\theta,u_z)(r,-\theta,z,t),
\]
and the force transforms identically. Pressure reflects in its argument.
For an axisymmetric field this reverses exactly its swirl; meridional
components remain fixed. The entire momentum tensor is mapped to
$O\mathcal T O^T$, so its radial--tangential pair transforms by
$(T_{r\theta},T_{rz})\mapsto(-T_{r\theta},T_{rz})$ after angular
averaging. Vorticity transforms as
$\omega^\sharp(x,t)=\det(O)\,O\omega(Ox,t)$, obtained by applying
the same derivative identities to the Levi-Civita formula for curl.
The map is an involution, preserves the kinetic-energy integral
(by $|\det O|=1$), compact support, smoothness, and blowup magnitude
along the reflected path. It preserves zero initial data when present.
It proves transport of any established axisymmetric signed swirl
divergence to the opposite sign; it does not itself prove an endpoint.

The source five-moment correction, auxiliary mean inverse, and infinite
summation are subsequent source operations. Their full estimates are
not re-proved here. The earlier draft formulas for them are withdrawn:
in particular the fifth moment row is
$\int(2RG\gamma_d-RV\Delta v)\,\dd R=-J_z$, without an extra outer
$R$, and its reserved profile in a $Q$ chart includes $(Q/q)^A$.
The finite map proved above retains the two moments and all other
residual terms, so it establishes neither a third shooting return nor
an infinite modified viscous sequence.


\section{The released result and its actual signed path}
\label{nsprop:source}
The source in this section is the supplied 166-page edition of
\emph{Finite Time Blowup for Navier--Stokes}, whose SHA-256 is
\begin{center}\footnotesize
\texttt{0e779481c4da40bd28d1e642e1d8ca57447d129610df28dfa5a11e9af8ae228f}.
\end{center}
Its Theorem 1.1, on page 1, states the following result for the standard
three-dimensional incompressible Navier--Stokes equation. For every
fixed physical viscosity $\nu_{\rm NS}>0$, the source constructs
\[
 f\in C^\infty_c(\mathbb R^3\times(0,\infty);\mathbb R^3),
 \qquad u,p\in C^\infty(\mathbb R^3\times[0,1)),
\]
with a fixed compact spatial support $K$ for $u$ and $p$, such that
\begin{equation}
 \begin{gathered}
 \partial_tu+(u\cdot\nabla)u-\nu_{\rm NS}\Delta u+\nabla p=f,
 \quad \nabla\cdot u=0,\quad u(\cdot,0)=0,\\
 \quad \sup_{t<1}\|u(t)\|_2<\infty,\quad
 \limsup_{t\uparrow1}\|u(t)\|_\infty=\infty .
 \end{gathered}
 \label{nsprop:source_theorem}
\end{equation}
It also excludes a global smooth solution with the same force and datum
and uniformly bounded kinetic energy. Its Corollary 10.6, pages 125--126,
gives a periodic construction on $\mathbb R^3/\mathbb Z^3$ with smooth
time-compact force, zero initial datum and no global smooth continuation.
The source identifies these with alternatives (C) and (D) in the
referenced Millennium problem statement. The force is nonzero, as the
source explicitly notes on page 124. These are statements of the cited
source theorem; the finite algebra checks in this lane do not constitute
an independent verification of the complete 166-page proof.

The source proves more than an unsigned norm divergence. In its
viscosity-one construction, equations (10.20)--(10.21), page 124, give
\begin{equation}
 \begin{split}
 A&=\tfrac12+h,\qquad 0<h<1/100,\qquad X_{\rm in}>0,\quad e_0>0,\\
 x_\tau&=(\sqrt{2X_{\rm in}\tau},0,0),\qquad t_\tau=1-\tau,\\
 u_\theta(x_\tau,t_\tau)
   &=\tau^{-A}(e_0+\epsilon(\tau)),\qquad
      |\epsilon(\tau)|\le C\tau^{2h}\quad(0<\tau<\tau_0).
 \end{split}\label{nsprop:positive_path}
\end{equation}
Here $C,\tau_0$ are finite constants supplied by the source asymptotic;
$e_0$ and $X_{\rm in}$ retain their source values. If $C>0$, choose
$0<\tau_1\le\tau_0$ so that $C\tau_1^{2h}\le e_0/2$; if $C=0$,
take $\tau_1=\tau_0$. It follows directly that
\[
 u_\theta(x_\tau,t_\tau)\ge \frac{e_0}{2}\tau^{-A}
       \longrightarrow+\infty\qquad(0<\tau<\tau_1).
\]
The observation point moves to the origin. This assertion does not
replace that path by the fixed point $x=0$. On the path, $q=\tau>0$,
and at its cylindrical angle $\theta=0$ one has $e_\theta=e_2$.
The geometric domain $q>0$, the positive exponent $h$, and the signed
velocity value are three retained parts of the same explicit formula.

\section{Exact propagation at retained viscosities and orientations}
\label{nsprop:map}
For any smooth fields $u,p$ define the physical momentum residual
\[
 R_{\nu_b}(u,p)=\partial_tu+(u\cdot\nabla)u-\nu_b\Delta u+\nabla p.
\]
Let $\nu_b>0$ and $\nu_t>0$ be the base and target physical viscosities,
and retain
\[
 \rho=\sqrt{\nu_t/\nu_b},\qquad O^TO=OO^T=I,\qquad a\in\mathbb R^3.
\]
Define the affine spatial coordinate $y=O^T(x-a)/\rho$ and the fields
\begin{equation}
 u^\sharp(x,t)=\rho O u(y,t),\quad
 p^\sharp(x,t)=\rho^2p(y,t),\quad
 f^\sharp(x,t)=\rho O f(y,t).
 \label{nsprop:affine_map}
\end{equation}
No time coordinate is changed. These formulas give a bijection on smooth
field triples, with inverse
\[
 u(y,t)=\rho^{-1}O^Tu^\sharp(a+\rho Oy,t),\quad
 p(y,t)=\rho^{-2}p^\sharp(a+\rho Oy,t),\quad
 f(y,t)=\rho^{-1}O^Tf^\sharp(a+\rho Oy,t).
\]
For the Jacobian convention $(Du)_{ij}=\partial_j u_i$, differentiation
gives, without dropping an orientation or viscosity,
\begin{align*}
 Du^\sharp&=O(Du)(y,t)O^T,&
 \nabla\cdot u^\sharp&=(\nabla\cdot u)(y,t),\\
 \partial_tu^\sharp&=\rho O\partial_tu(y,t),&
 (u^\sharp\cdot\nabla_x)u^\sharp&=\rho O((u\cdot\nabla_y)u)(y,t),\\
 \Delta_xu^\sharp&=\rho^{-1}O\Delta_yu(y,t),&
 \nabla_xp^\sharp&=\rho O\nabla_yp(y,t).
\end{align*}
The nonlinear identity follows by multiplying $Du^\sharp$ by
$u^\sharp$ and using $O^TO=I$. The Laplacian identity follows by
contracting the two spatial derivative indices and using orthogonality.
Since $\nu_t=\rho^2\nu_b$, these identities prove
\begin{equation}
 R_{\nu_t}(u^\sharp,p^\sharp)-f^\sharp
       =\rho O\bigl(R_{\nu_b}(u,p)-f\bigr)(y,t).
 \label{nsprop:residual_map}
\end{equation}
In particular the map preserves the exact equation as well as each
nonzero residual with its exact factor.

For vorticity, the cross-product identity
$(Ov)\times(Ow)=\det(O)O(v\times w)$ gives
\begin{equation}
 \omega^\sharp(x,t)=\det(O)O\omega(y,t).
 \label{nsprop:vorticity}
\end{equation}
One proof of this identity takes its scalar product with $Oz$:
both sides give $\det(O)\det[v,w,z]$ for every $z$. Applying it to
the contracted first derivatives proves the vorticity formula; the
factor $\rho$ in velocity cancels the inverse factor in a derivative.

Spatial support becomes $a+\rho OK$. The Jacobian has absolute
determinant $\rho^3$, for both signs of $\det O$. Hence, for
$1\le p<\infty$,
\begin{equation}
 \|u^\sharp(t)\|_p=\rho^{1+3/p}\|u(t)\|_p,\qquad
 \|u^\sharp(t)\|_\infty=\rho\|u(t)\|_\infty,\qquad
 \|u^\sharp(t)\|_2^2=\rho^5\|u(t)\|_2^2 .
 \label{nsprop:norms}
\end{equation}
The Frobenius norm of $O(Du)O^T$ equals that of $Du$. Consequently,
with integrals over the same time interval,
\[
 \nu_t\int\|\nabla u^\sharp\|_2^2\,dt
     =\rho^5\nu_b\int\|\nabla u\|_2^2\,dt .
\]
For arbitrary directions $v_1,\ldots,v_j$, the exact force derivative is
\[
 D_x^j\partial_t^m f^\sharp[v_1,\ldots,v_j](x,t)
 =\rho^{1-j}O
   (D_y^j\partial_t^m f)[O^Tv_1,\ldots,O^Tv_j](y,t).
\]
Thus smoothness, compact space-time support and vanishing of every
derivative at a terminal point are preserved. The initial datum remains
zero. A global smooth bounded-energy competitor for the transformed
datum and force would, under the displayed inverse, be a competitor
for the base datum and force, because its energy is multiplied by the
fixed factor $\rho^{-5}$. This proves propagation of the source
noncontinuation conclusion through the map.

\section{The source positive branch and the exact negative branch}
\label{nsprop:branches}
Apply the map to the source field in \eqref{nsprop:positive_path}.
Here $\nu_b=1$ specifies the source object being used; physical target
viscosity $\nu_t$ remains arbitrary and $\rho=\sqrt{\nu_t}$ is retained.
Take $a=0$ and the two particular orthogonal matrices
\[
 O_+=I,\qquad O_-=\operatorname{diag}(1,-1,1).
\]
Both fix $e_1$, so they give the same spatial observation path
$x^\sharp_\tau=\rho x_\tau$ and the same time $1-\tau$.
Since $O_+e_2=e_2$ and $O_-e_2=-e_2$, their exact signed values are
\begin{equation}
 u^\sharp_{\pm,\theta}(\rho x_\tau,1-\tau)
     =\pm\rho\,\tau^{-A}(e_0+\epsilon(\tau))
        \longrightarrow\begin{cases}+\infty,&+,\\-\infty,&-.\end{cases}
 \label{nsprop:two_paths}
\end{equation}
All support, energy, force and viscosity statements in the preceding
section apply to both triples. For the negative branch the force is
$\rho O_-f(O_-x/\rho,t)$. Equality of this force with the positive
branch force is not needed and is not asserted.

For a general, possibly nonaxisymmetric field the reflection alone has
the exact cylindrical component formula
\[
 (u^r,u^\theta,u^z)^\sharp(r,\theta,z,t)
    =(u^r,-u^\theta,u^z)(r,-\theta,z,t).
\]
Indeed $O_-e_r(-\theta)=e_r(\theta)$,
$O_-e_\theta(-\theta)=-e_\theta(\theta)$ and $O_-e_z=e_z$.
This proves the formula even in the presence of all oscillatory
corrections. For the axisymmetric background the argument $-\theta$
can be dropped because its components are independent of $\theta$.
For a velocity covariance tensor the exact Cartesian transformation is
$T^\sharp(x,t)=\rho^2 O T(y,t)O^T$: expand the outer product of
$\rho Ow$ with itself and use linearity of the retained average.
For its angularly averaged symmetric tensor, the reflection basis
calculation at $\rho=1$ gives
\[
 (T_{r\theta},T_{rz})^\sharp=(-T_{r\theta},T_{rz}).
\]
Under the additional viscosity dilation its components have factor
$\rho^2$ at the corresponding points. Thus the two covariance columns
and their target transform together. Positive coefficients in a source
representation $T=c_1v_1+c_2v_2$, $c_i>0$, remain the same coefficients
of the transformed columns; no change of coefficient sign is required
to realize the negative-swirl branch.

The source's periodic construction also has this reflection: $O_-$
preserves $\mathbb Z^3$, so $[x]\mapsto[O_-x]$ is well defined on
$\mathbb R^3/\mathbb Z^3$. The derivative identities apply to periodic
lifts. This gives both signed orientations of the source Corollary 10.6.
Arbitrary $O\in O(3)$ need not preserve this particular lattice; only
the lattice-preserving reflection just specified is used for this
periodic assertion.

\section{Propagation into the coupled calculation}
\label{nsprop:programme}
The preceding sections use the released whole-space and periodic
constructions at the precise places where those source results apply.
The original coupled Boussinesq work is due to Levent Alp\"oge and
Tristan Buckmaster. This lane's retained coupled system, its finite
controlled stages and its residual maps remain the separate objects
defined in the preceding parts of this reader.

For every compact candidate $(u_*+w_y,p_*+\pi_y)$ in the covariance
part, \eqref{nsprop:residual_map} propagates its full computed residual.
Linearity of the transformation and the exact nonlinear identity give
in particular
\[
 \Delta R_y^\sharp(x,t)=\rho O\Delta R_y(y,t).
\]
Thus the finite candidate's cutoff forces, curl errors and covariance
remainders all have an exact transformed counterpart. A nonzero
remainder is carried to a nonzero remainder, because the map is
invertible. The source's residual flatness at the singular point belongs
to its completed local source construction; its extended smooth force
need not vanish on the entire time-one slice. The next section calculates the
surviving moments of the actual finite candidate, so that propagation
continues through its computed terms rather than assuming that its
different coupled coefficients inherit the source's infinite estimates.

The earlier suggestion of a negative blowout imposes no condition on
the programme. Formula \eqref{nsprop:positive_path} is the source-selected
sign; formula \eqref{nsprop:two_paths} proves the exact relation of its
two orientations. Both formulas retain a positive physical viscosity
and the source coordinate domain $q>0$.


\section{The two moments of the actual compact curl increment}
\label{nsmom:setup}
The objects in this calculation are the finite smooth source state
$(u_*,p_*)$ and the actual added curl field $v=w_y$ with its selected
pressure $\pi=\pi_y$ from the signed covariance bridge. Their physical
residual increment is
\begin{equation}
 \Delta R=\partial_t v+(u_*\cdot\nabla)v+(v\cdot\nabla)u_*
 +(v\cdot\nabla)v-\nu_{\rm NS}\Delta v+\nabla\pi .
 \label{nsmom:increment}
\end{equation}
The physical viscosity $\nu_{\rm NS}$ is unchanged. Both velocities are
divergence free: the source state uses its actual divergence-free
background and curl increments, and $v$ is the curl of its specified
compact potential. No residual equation for $u_*$ is required in
\eqref{nsmom:increment}; its residual has been subtracted exactly.
All radial supports involving $v$ or $\pi$ lie in a fixed compact
subinterval of $(0,\infty)$ on the compact space--time patch at issue.

Here the source extended field is indispensable. Its absolute auxiliary
coordinate is $Y\in\mathbb T^2=\mathbb R^2/\mathbb Z^2$, and physical
evaluation uses precisely source (6.2)--(6.4):
\begin{gather}
 J_g=\begin{pmatrix}3&1\\1&5\end{pmatrix},\quad
 \Lambda_g=4-\sqrt2,\quad T_g=4+\sqrt2,\quad b_g=\sqrt2-1,\nonumber\\
 \mathbf v_r=(1,-b_g),\quad\mathbf v_t=(b_g,1),\quad
 \rho_g=\frac{\log\Lambda_g}{\log T_g},\quad
 \kappa_s=10^{-5},\quad d_r=2((1+h)\rho_g-h\kappa_s),\nonumber\\
 Y(r,t)=\mathbf v_r r^{d_r}+\mathbf v_t t\pmod{\mathbb Z^2},\qquad
 \mathscr D_r=\partial_r+d_r r^{d_r-1}\mathbf v_r\cdot\partial_Y,
 \quad\mathscr D_t=\partial_t+\mathbf v_t\cdot\partial_Y.
 \label{nsmom:phase}
\end{gather}
These are the source's original quantities, not the coupled physical
frequency or diffusivity. Angular and axial physical derivatives are
$\partial_\theta$ and $\partial_z$, because $Y(r,t)$ depends on neither
variable. Every field below is first differentiated using these lifted
physical derivatives. A tilde is suppressed for this extended
representation. Evaluation at \eqref{nsmom:phase} recovers the physical
field and its physical derivatives by the chain rule.

The incompressibility needed below holds at every independent $Y$.
Indeed the actual finite source construction retains the form
$u_* =\operatorname{curl}_{\mathscr D}\mathcal A_*+
B_*e_\theta$ with $\partial_\theta B_*=0$, and the added field is
$v=\operatorname{curl}_{\mathscr D}\mathcal A_y$. The retained source
potentials here include its finite mean and wave curl increments.
The lifted derivatives in \eqref{nsmom:phase}, together with
$\partial_\theta$ and $\partial_z$, commute, and
$\mathscr D_r r=1$. Substitution in the cylindrical curl and
divergence formulas therefore gives
$\operatorname{div}_{\mathscr D}\operatorname{curl}_{\mathscr D}
\mathcal A=0$ identically in $(r,\theta,z,t,Y)$: each mixed derivative
cancels its opposite, including the derivative of $r^{-1}$.
Also $\operatorname{div}_{\mathscr D}(B_*e_\theta)
=r^{-1}\partial_\theta B_*=0$.
Thus extended incompressibility is proved from the actual retained
potential representation, rather than inferred from its restriction to
the physical phase map.

Use the normalized averages
\[
 \langle a\rangle_\theta=\frac1{2\pi}\int_0^{2\pi}a\,d\theta,
 \qquad \langle a\rangle_Y=\int_{\mathbb T^2}a\,dY,
 \qquad\langle a\rangle=\langle\langle a\rangle_\theta\rangle_Y,
 \qquad\int_{\mathbb T^2}1\,dY=1.
\]
They act on cylindrical components before physical evaluation, exactly
as source (3.7)--(3.9). All claimed integral identities below concern
these specified averages; the exact relation to evaluated moments is
also proved below.

\subsection{The actual harmonic curl has zero angular mean}
\label{nsmom:zero_mean}
For each retained source harmonic, its extended potential and pressure
have cylindrical components of the form
\[
 C_j(r,z,t,Y)e^{i(n_j\theta+\psi_j(r,z,t,Y))},\qquad
 P_j(r,z,t,Y)e^{i(n_j\theta+\psi_j(r,z,t,Y))},\qquad
 n_j=k_{\beta_j}m_jp_j\in\mathbb Z\setminus\{0\}.
\]
The coefficients include the actual cutoff $\chi$, frame $B$, map $L$,
coupled solution $y=(\Theta,\Omega)^T$, band powers and pressure formula.
Their cylindrical components are independent of $\theta$. In particular,
no sign of $\Theta$, $\Omega$ or the real part of a harmonic is prescribed.
For a vector potential $\mathcal A$ the exact cylindrical curl is
\[
 (\operatorname{curl}\mathcal A)_r=r^{-1}\partial_\theta\mathcal A_z
                         -\partial_z\mathcal A_\theta,
 \quad
 (\operatorname{curl}\mathcal A)_\theta=\partial_z\mathcal A_r
                         -\mathscr D_r\mathcal A_z,
 \quad
 (\operatorname{curl}\mathcal A)_z=(\mathscr D_r+r^{-1})\mathcal A_\theta
                         -r^{-1}\partial_\theta\mathcal A_r.
\]
These formulas include the derivatives of the cylindrical basis.
Each displayed operation retains the same factor $e^{in_j\theta}$:
the coefficients of $\mathscr D_r$ are independent of $\theta$, and
$\partial_\theta$ multiplies that factor by $in_j$. Integration of the
factor over the angular circle gives zero. Taking real parts and the
actual finite harmonic sum proves
\begin{equation}
 \langle v_r\rangle_\theta=\langle v_\theta\rangle_\theta
 =\langle v_z\rangle_\theta=\langle\pi\rangle_\theta=0
 \quad\hbox{at every }(r,z,t,Y).
 \label{nsmom:mean_zero}
\end{equation}
This includes every curl and cutoff correction, not merely the
principal amplitude. The same identities hold after physical
evaluation because $Y(r,t)$ is independent of $\theta$.

For every smooth periodic coefficient $a$, integration of its auxiliary
derivatives gives
\begin{equation}
 \langle\mathscr D_r a\rangle=\partial_r\langle a\rangle,
 \quad\langle\mathscr D_t a\rangle=\partial_t\langle a\rangle,
 \quad\langle\partial_z a\rangle=\partial_z\langle a\rangle,
 \quad\langle\partial_\theta a\rangle=0.
 \label{nsmom:commute}
\end{equation}
There is no omitted variable coefficient in this assertion:
$d_r r^{d_r-1}\mathbf v_r$ is independent of $Y$. Applying the first
identity again proves
$\langle\mathscr D_r^2a\rangle=\partial_r^2\langle a\rangle$.
Explicitly, if $c(r)=d_r r^{d_r-1}\mathbf v_r$, then
\[
 \mathscr D_r^2 a=\partial_r^2a+2c\cdot\partial_Y\partial_r a
 +c'\cdot\partial_Ya+(c\cdot\partial_Y)^2a;
\]
each of the three additional terms has zero auxiliary average.

\subsection{Exact removal of the linear and mean-flow terms}
\label{nsmom:linear}
Put
\[
 \bar u_*:=\langle u_*\rangle_\theta,
 \qquad w_*:=u_*-\bar u_*.
\]
The field $\bar u_*$ may retain auxiliary dependence; none is discarded.
Since it is independent of $\theta$, \eqref{nsmom:mean_zero} gives
\[
 \langle\bar u_*\otimes v+v\otimes\bar u_*\rangle=0.
\]
Define the actual symmetric cross-and-quadratic tensor
\begin{equation}
 S:=\langle w_*\otimes v+v\otimes w_*+v\otimes v\rangle.
 \label{nsmom:stress}
\end{equation}
Products of all original source waves with the added curl field are
retained in this formula, including old and new curl remainders.

For two divergence-free fields $a,b$, Cartesian component expansion
gives
\[
 \nabla\cdot(a\otimes b+b\otimes a)
 =(b\cdot\nabla)a+(a\cdot\nabla)b,
 \qquad\nabla\cdot(v\otimes v)=(v\cdot\nabla)v.
\]
For example
$\sum_j\partial_j(a_i b_j)=\sum_jb_j\partial_ja_i+
a_i\sum_j\partial_jb_j$; the last term is zero. The lifted derivatives
obey the same product rule, so the identities hold on the extended
domain before evaluation. The three advection terms in
\eqref{nsmom:increment} are therefore exactly the divergence of
$u_*\otimes v+v\otimes u_*+v\otimes v$.

For completeness, set
$\mathscr L_0=\mathscr D_r^2+r^{-1}\mathscr D_r+
r^{-2}\partial_\theta^2+\partial_z^2$.
The vector Laplacian and pressure gradient in the two components are
\[
 (\Delta v)_\theta=(\mathscr L_0-r^{-2})v_\theta
                       +2r^{-2}\partial_\theta v_r,
 \qquad(\Delta v)_z=\mathscr L_0v_z,
 \qquad(\nabla\pi)_\theta=r^{-1}\partial_\theta\pi,
 \qquad(\nabla\pi)_z=\partial_z\pi.
\]
Equations \eqref{nsmom:mean_zero}--\eqref{nsmom:commute} prove separately
\begin{equation}
 \langle\mathscr D_t v\rangle_{\rm tan}=0,
 \qquad -\nu_{\rm NS}\langle\Delta v\rangle_{\rm tan}=0,
 \qquad\langle\nabla\pi\rangle_{\rm tan}=0.
 \label{nsmom:linear_zero}
\end{equation}
Thus viscosity and pressure were included before averaging; their
vanishing here is a proved consequence of the actual carriers.

If $T$ is symmetric, differentiation of
$e_r=(\cos\theta,\sin\theta,0)$,
$e_\theta=(-\sin\theta,\cos\theta,0)$ gives
\begin{align*}
 (\nabla\cdot T)_\theta
   &=(\mathscr D_r+2/r)T_{r\theta}
       +r^{-1}\partial_\theta T_{\theta\theta}+\partial_zT_{z\theta},\\
 (\nabla\cdot T)_z
   &=(\mathscr D_r+1/r)T_{rz}
       +r^{-1}\partial_\theta T_{\theta z}+\partial_zT_{zz}.
\end{align*}
Indeed $\partial_\theta e_r=e_\theta$ and
$\partial_\theta e_\theta=-e_r$ supply the two radial connection
terms in the first line and the single radial connection term in the
second. Averaging these formulas and using
\eqref{nsmom:commute}, \eqref{nsmom:stress}, and \eqref{nsmom:linear_zero}
proves the exact residual identities
\begin{equation}
 F_2:=\langle\Delta R_\theta\rangle
    =(\partial_r+2/r)S_{r\theta}+\partial_zS_{z\theta},
 \qquad
 F_1:=\langle\Delta R_z\rangle
    =(\partial_r+1/r)S_{rz}+\partial_zS_{zz}.
 \label{nsmom:F}
\end{equation}
In particular every mean-flow cross term vanishes as the divergence
of its zero averaged tensor, even when the mean depends on $Y$.

\subsection{The surviving axial fluxes and their signs}
\label{nsmom:moments}
Define the two moments with the original physical radial measure:
\begin{equation}
 M_2(z,t)=\int_0^\infty r^2F_2(r,z,t)\,dr,
 \qquad M_1(z,t)=\int_0^\infty rF_1(r,z,t)\,dr.
 \label{nsmom:Mdef}
\end{equation}
Their exact values are
\begin{align}
 M_2&=\partial_z J_2, &
 J_2&=\int_0^\infty r^2
    \langle w_{*,z}v_\theta+v_z w_{*,\theta}+v_zv_\theta\rangle\,dr,
 \label{nsmom:M2}\\
 M_1&=\partial_z J_1, &
 J_1&=\int_0^\infty r
    \langle2w_{*,z}v_z+v_z^2\rangle\,dr.
 \label{nsmom:M1}
\end{align}
To prove them, multiply \eqref{nsmom:F} by the respective weight:
$r^2(\partial_r+2/r)S_{r\theta}=\partial_r(r^2S_{r\theta})$
and $r(\partial_r+1/r)S_{rz}=\partial_r(rS_{rz})$.
The weighted boundary values at both ends of $(0,\infty)$ are zero,
because each entry of $S$ contains at least one compactly supported
factor $v$. Differentiation under the radial integral is valid by the
same compact support and smoothness. This proves
\eqref{nsmom:M2}--\eqref{nsmom:M1} including their signs and factors.

There is no pointwise cancellation of the axial derivatives in these
formulas. Since the added field has compact axial support, $J_1$ and
$J_2$ have compact axial support and
\begin{equation}
 \int_{\mathbb R}M_e(z,t)\,dz=0\quad(e=1,2).
 \label{nsmom:global_zero}
\end{equation}
This follows by evaluating the boundary values of $J_e$ at the two
axial ends. It does not assert $M_e(z,t)=0$ at an individual axial
point. The compact radial primitive in the covariance bridge must
therefore retain its $b_eM_e$ terms with these actual $M_e$.

At fixed source state, replacing the added field and pressure by their
negatives sends the old-wave cross terms in \eqref{nsmom:M2}--\eqref{nsmom:M1}
to their negatives and leaves the quadratic terms unchanged. Thus
negating an amplitude is not a sign reversal of either full moment.
Under the full physical reflection
$O=\operatorname{diag}(1,-1,1)$, applied to the source state and added
field together as $a^\sharp(x,t)=Oa(Ox,t)$, the cylindrical radial and
axial components retain their signs and the angular components reverse
sign at $(r,-\theta,z,t)$. Changing the angular integration variable
therefore proves
\[
 J_2^\sharp=-J_2,\quad M_2^\sharp=-M_2,
 \qquad J_1^\sharp=J_1,\quad M_1^\sharp=M_1.
\]

An explicit transverse-curl example shows why compactness and nonzero
angular frequency alone do not force either moment to vanish. This
example tests precisely those structural properties, and is not a
replacement for the retained source state. Let $H$ be any nonzero real
$C_c^\infty((r_-,r_+))$ function with $0<r_-<r_+$, let
$g\in C_c^\infty(\mathbb R)$ be nonconstant, let $n$ be a nonzero
integer, and put $f=rH'$. Use zero old state and the real potential
\[
 \mathcal A_r=fg\sin(n\theta),\qquad
 \mathcal A_\theta=0,\qquad
 \mathcal A_z=Hg\cos(n\theta).
\]
It has the required transverse representation:
$\phi=n\theta$, $\xi=(n/r)e_\theta$,
$C=(-ifg,0,Hg)$ in cylindrical components,
$a=i\xi\times C$, and
$C=i\xi\times a/|\xi|^2$ because $\xi\cdot C=0$.
The curl is exactly
\[
 v_r=-\frac{nHg}{r}\sin(n\theta),\qquad
 v_\theta=fg'\sin(n\theta)-H'g\cos(n\theta),\qquad
 v_z=-\frac{nfg}{r}\cos(n\theta).
\]
Consequently
\[
 M_2=ngg'\int_0^\infty r^2(H')^2\,dr,
 \qquad
 M_1=n^2gg'\int_0^\infty r(H')^2\,dr.
\]
For example $g(z)=\exp(-1/(1-z^2))$ on $|z|<1$, extended by zero,
has $g(z)g'(z)\ne0$ at every $0<|z|<1$. Both integrals are strictly
positive: a compactly supported smooth function with $H'=0$ everywhere
would be constant and hence zero. Thus both moments can be nonzero,
with their displayed signed values, within the exact transverse-curl
architecture.

\subsection{Bounds in the actual fields and original chart units}
\label{nsmom:bounds}
At fixed $(z,t)$ set
\[
 \|a\|_e^2=\int_0^\infty r^e\langle|a|^2\rangle\,dr,
 \qquad e=1,2.
\]
When an old-field norm is used below it can be restricted to the
enclosing radial support of $v$ and $\partial_zv$. This avoids imposing
any decay condition on an unused exterior part of the old field.
The following component bounds contain only actual fields:
\begin{align}
 |J_2|\le{}&\|w_{*,z}\|_2\|v_\theta\|_2
 +\|w_{*,\theta}\|_2\|v_z\|_2+\|v_z\|_2\|v_\theta\|_2,
 \label{nsmom:J2bound}\\
 |J_1|\le{}&2\|w_{*,z}\|_1\|v_z\|_1+\|v_z\|_1^2,
 \label{nsmom:J1bound}\\
 |M_2|\le{}&\|\partial_z w_{*,z}\|_2\|v_\theta\|_2
 +\|w_{*,z}\|_2\|\partial_zv_\theta\|_2
 +\|\partial_zv_z\|_2\|w_{*,\theta}\|_2\nonumber\\
 &+\|v_z\|_2\|\partial_z w_{*,\theta}\|_2
 +\|\partial_zv_z\|_2\|v_\theta\|_2
 +\|v_z\|_2\|\partial_zv_\theta\|_2,
 \label{nsmom:M2bound}\\
 |M_1|\le{}&2\bigl(\|\partial_z w_{*,z}\|_1\|v_z\|_1
 +\|w_{*,z}\|_1\|\partial_zv_z\|_1
 +\|v_z\|_1\|\partial_zv_z\|_1\bigr).
 \label{nsmom:M1bound}
\end{align}
Each is the Cauchy--Schwarz inequality on the product measure
$r^e\,dr\,d\theta/(2\pi)\,dY$, after applying the product rule in
\eqref{nsmom:M2}--\eqref{nsmom:M1}. For example, the derivative of
$2w_{*,z}v_z+v_z^2$ is
$2(\partial_z w_{*,z})v_z+2w_{*,z}\partial_zv_z+
2v_z\partial_zv_z$, proving every coefficient in
\eqref{nsmom:M1bound}.
For a $(z,t)$ multi-index $I$, every further derivative is exactly
obtained by replacing each product $ab$ in $J_e$ by
\[
 \partial^{I+(1,0)}(ab)
 =\sum_{K\le I+(1,0)}\binom{I+(1,0)}{K}
       (\partial^Ka)(\partial^{I+(1,0)-K}b)
\]
inside its original radial integral. This identity follows inductively
from the one-derivative product rule, and gives the corresponding
finite-order moment bound by applying the same Cauchy--Schwarz
inequality separately to every displayed summand.

The field $v$ in these bounds is not an abstract replacement amplitude.
For each actual band and harmonic write
\begin{gather}
 t_j=\chi_jB_jL_jy_j,\qquad
 c_j=\frac{i\,n_{\Phi,j}\times t_j}
                  {k_{\beta_j}m_j|n_{\Phi,j}|^2},\qquad
 b_j=t_j+\operatorname{curl}_{*,\mathrm{coeff}}c_j,\nonumber\\
 v=\sum_j\operatorname{Re}
       \bigl(Q_{\beta_j}^{-A}b_j e^{ik_{\beta_j}m_j\Phi_j}\bigr),
 \qquad A=\tfrac12+h,\quad D=\tfrac12-h.
 \label{nsmom:actual_amplitude}
\end{gather}
Here $\operatorname{curl}_{*,\mathrm{coeff}}$ differentiates the chart
coefficient, including the cutoff, auxiliary chain rule and cylindrical
basis, but not the displayed exponential. This is exactly the
remainder of the physical potential
$Q_{\beta_j}^{1/2-A}c_j e^{ik_{\beta_j}m_j\Phi_j}$: physical curl
introduces the factor $Q_{\beta_j}^{-1/2}$, and differentiating the
exponential produces $t_j$ by the transverse triple-product identity.
Explicitly, on the absolute auxiliary torus in a fixed band,
\[
 \mathscr D_R=\partial_R+
 \sqrt Q\,d_r r^{d_r-1}\mathbf v_r\cdot\partial_Y,
 \qquad r=\sqrt Q R,\qquad \mathscr D_Z^*=\varepsilon\partial_Z.
\]
For the actual coefficient $c=(c_r,c_\theta,c_z)$, which has
$\partial_\theta c=0$, the operator is exactly
\[
 \operatorname{curl}_{*,\mathrm{coeff}}c
 =\bigl(-\varepsilon\partial_Zc_\theta,
       \varepsilon\partial_Zc_r-\mathscr D_Rc_z,
       (\mathscr D_R+R^{-1})c_\theta\bigr).
\]
Coefficients originally on a band auxiliary torus are pulled back to
the absolute torus before applying this formula; this retains the
integer-covering chain derivative in $\partial_Yc$.
In the same fixed chart the exact axial derivative is
\begin{align}
 \partial_zv
 &=\sum_j\operatorname{Re}\left[
 Q_{\beta_j}^{-A-D}
 \bigl(\partial_{Z_j}b_j+
   ik_{\beta_j}m_j(\partial_{Z_j}\Phi_j)b_j\bigr)
 e^{ik_{\beta_j}m_j\Phi_j}\right].
 \label{nsmom:actual_derivative}
\end{align}
The band power is fixed during chart differentiation and
$\partial_z=Q_{\beta_j}^{-D}\partial_{Z_j}$; auxiliary phase evaluation
adds no axial derivative because of \eqref{nsmom:phase}. Thus
\eqref{nsmom:actual_amplitude}--\eqref{nsmom:actual_derivative} substitute
the actual retained pair and all its curl terms into the bounds.
For example, writing
$\|b\|_{e,\beta}^2=\int R^e\langle|b|^2\rangle\,dR$ in that band,
the triangle inequality and $|\operatorname{Re}z|\le|z|$ give
\begin{align}
 \|v\|_e
 &\le\sum_j Q_{\beta_j}^{(e+1)/4-A}\|b_j\|_{e,\beta_j},
 \label{nsmom:scale_bound}\\
 \|\partial_zv\|_e
 &\le\sum_j Q_{\beta_j}^{(e+1)/4-A-D}
 \|\partial_{Z_j}b_j+
 ik_{\beta_j}m_j(\partial_{Z_j}\Phi_j)b_j\|_{e,\beta_j}.
 \label{nsmom:scale_derivative_bound}
\end{align}
The exponent $(e+1)/4$ is exactly the square root of the radial measure
factor $r^e\,dr=Q^{(e+1)/2}R^e\,dR$. Distinct bands and harmonics
remain distinct in these sums. No limiting estimate as $q\downarrow0$
has been inferred from the finite compact-patch norms.

For a single fixed chart, put
$F_e=Q^{-2A-1/2}E_e^*$ as in the covariance bridge and
$S^*=Q^{2A}S$. Direct change of radial variable gives
\[
 M_e=Q^{e/2-2A}M_e^*,\qquad
 M_e^*=\int R^eE_e^*\,dR,
 \qquad
 M_e^*=\varepsilon\partial_Z\int R^e S^*_{z,\tau_e}\,dR,
 \quad\varepsilon=Q^h,
\]
where $\tau_2=\theta$ and $\tau_1=z$. Indeed the axial flux has
factor $Q^{(e+1)/2-2A}$, its axial derivative contributes $Q^{-D}$,
and dividing by $Q^{e/2-2A}$ leaves $Q^{1/2-D}=Q^h$.
This proves the precise relation to the source chart moment without
altering its axial anisotropy.

\subsection{Exact relation to moments after phase evaluation}
\label{nsmom:evaluated}
Let
\[
 \mathcal S=\langle w_*\otimes v+v\otimes w_*+v\otimes v\rangle_\theta,
 \qquad \mathcal S^\circ=\mathcal S-\langle\mathcal S\rangle_Y.
\]
Thus $\langle\mathcal S\rangle_Y=S$ and
$\langle\mathcal S^\circ\rangle_Y=0$. For the evaluated physical
fields, denote the angular-only residual moments by $M_e^{\rm ev}$.
The proof of \eqref{nsmom:mean_zero} holds after evaluation. The
physical divergence, vector Laplacian and gradient identities therefore
prove, without any auxiliary averaging,
\[
 M_e^{\rm ev}
   =\partial_z\int_0^\infty r^e
        \mathcal S_{z,\tau_e}(r,z,t,Y(r,t))\,dr.
\]
The physical radial boundary terms still vanish by compact support;
the chain rule is included in the physical radial derivative before
that integration. Subtracting \eqref{nsmom:M2}--\eqref{nsmom:M1} proves
the exact comparison
\begin{equation}
 M_e^{\rm ev}-M_e
 =\partial_z\int_0^\infty r^e
     \mathcal S^\circ_{z,\tau_e}(r,z,t,Y(r,t))\,dr.
 \label{nsmom:evaluation_difference}
\end{equation}
For example its absolute value is bounded by
$\int r^e\sup_Y|\partial_z\mathcal S^\circ_{z,\tau_e}|\,dr$,
because $Y(r,t)$ is independent of $z$. Zero auxiliary Haar average
does not make the right side of \eqref{nsmom:evaluation_difference}
zero at the physical phase. The displayed map supplies precisely the
remaining auxiliary-oscillation flux; it is not silently identified
with the source fast mean.

\clearpage\part{Original third-growth derivatives}
\section{Time derivatives of the actual signed third growth}
\label{tgdc:section}

This calculation is on the actual interval already constructed in
\emph{The actual third growth target on the moving second hold} and
\emph{Exact third growth on the moving second holding state}. It gives
quantitative derivatives of that trajectory for the source-frame map.
All derivatives in this section are with respect to coupled physical
time $t_c$, and $\tau=t_c-t_b$ denotes elapsed coupled time. The notation
does not use the cylindrical spatial radius for elapsed time.

\subsection{Original data, evolving parent, and proved interval}

Retain the source parameters
\begin{equation}\label{tgdc:original}
\begin{gathered}
 \nu=\sqrt{A_0},\qquad
 \mu=\lambda_0=\lceil\sqrt{A_0}/2\rceil,\qquad
 K_j=\mu\widehat\lambda_j,\\
 S_3=\frac{A_0}{\mu}\widehat\lambda_3^{-k_3-6},\qquad
 P_3^*=\frac{A_0}{\mu}\widehat\lambda_3^{-7/8}=S_3e^{L_3}.
\end{gathered}
\end{equation}
The physical common diffusivity is $d>0$; it is distinct from $\nu$.
Use the actual second return $t_b$, its values $h_b,A_{1,b},P_{2,b}$,
and the original angles $s_*,s_2,s_3$. Put
\begin{equation}\label{tgdc:constants}
\begin{gathered}
 a=\sin s_2,\quad \chi_b=\sqrt{h_b^2+a^2},\quad
 \beta_b=\arctan(a/h_b),\quad
 b=\sin(s_3+\beta_b),\quad g_b=\cos(s_3+\beta_b),\\
 c=bh_b-ag_b=\chi_b\sin s_3,\qquad
 S_*=A_0\sin s_*,\quad C_*=A_0\cos s_*.
\end{gathered}
\end{equation}
Here $a,b,c,S_*,C_*>0$, and the letter $b$ is the original oriented
geometric constant, whereas $b_3$ below is an amplitude coefficient.
The moving parent is exactly
\begin{equation}\label{tgdc:parent}
\begin{aligned}
 \chi^2&=h^2+a^2,& \dot h&=a\Omega_1,&
 \dot\Omega_1&=aA_1/\chi-dK_1^2\Omega_1,\\
 \dot A_1&=-S_*\Omega_1-dK_1^2A_1,&
 \dot P_2&=-dK_2^2\chi^2P_2,&
 \dot g&=b\Omega_1,\\
 g&=g_b+(b/a)(h-h_b),&
 N_3&=b(A_1+C_*)+gS_*+cK_2P_2,&
 R_3&=g^2+b^2.
\end{aligned}
\end{equation}
The retained feedback is $\dot\alpha=-a^2\Omega_1/\chi^2$ and
$\Omega_2=0$. In particular the older temperature and first vorticity
are not frozen. The cancellation in the complete numerator is
\begin{equation}\label{tgdc:numerator}
 \dot N_3=-dE_3,\qquad
 E_3=bK_1^2A_1+cK_2^3\chi^2P_2.
\end{equation}

Write $N_i=N_3(t_b)$, $a_i=N_i/K_3$,
$b_i=K_3\sin s_3$, $\gamma_i=\sqrt{a_i b_i}$ and
$u_i=\sqrt{b_i/a_i}$. The exact new equations, including their
original signed seed, are
\begin{equation}\label{tgdc:actual_pair}
\begin{gathered}
 y=\binom{\Theta_3}{\Omega_3},\qquad
 \dot y=C_3y,\qquad
 C_3=\begin{pmatrix}-\lambda_3&a_3\\b_3&-\lambda_3\end{pmatrix},\\
 a_3=\frac{N_3}{K_3R_3},\qquad
 b_3=\frac{K_3c}{\chi},\qquad
 \lambda_3=dK_3^2R_3,\qquad
 y(t_b)=-S_3\binom1{u_i}.
\end{gathered}
\end{equation}
There are no new amplitude controls in this growth pair:
$c_\vartheta=c_\omega=0$. The activation of the compact wave is a
separate original factor $\eta(\Gamma_3\tau)$; it does not replace
\eqref{tgdc:actual_pair} by a forced coefficient trajectory.

For clarity, the following facts about this actual interval were proved
in the third-threshold construction. They identify the trajectory to
which the estimates below apply:
\begin{equation}\label{tgdc:proved_boxes}
\begin{gathered}
 I_3=[t_b,t_3^a],\quad t_3^a-t_b=\tau_a<14/\sigma_1,\quad
 P_3=-\Theta_3>0,\quad P_3(t_b)=S_3,\quad P_3(t_3^a)=P_3^*,\\
 u=\Omega_3/\Theta_3,\qquad
 \frac{u_i}{\sqrt2}\leq u\leq2u_i,\qquad
 \frac{a_i}4\leq a_3\leq a_i,\qquad
 \frac{b_i}2\leq b_3\leq b_i,\\
 1\leq R_3\leq\overline R:=\left(\frac{265}{256}\right)^2<2,
 \quad 1\leq\frac\chi{\chi_b}\leq\frac{259}{256},\quad
 0<dK_3^2\leq\frac{\gamma_i}{16\sqrt2},\\
 \frac{\gamma_i}{8\sqrt2}\leq (\log P_3)'\leq2\gamma_i,
 \quad 0<A_1\leq A_{1,b},\quad 0<P_2\leq P_{2,b},\quad
 0<\Omega_1<9\sigma_1,\\
 h_b\leq h\leq h_b+144a,\qquad
 g_b\leq g\leq g_b+144b,\qquad
 \tfrac12N_i\leq N_3\leq N_i,\qquad 1\leq\chi^2<10.
\end{gathered}
\end{equation}
The value of $d$ is precisely the previously selected range
$0<d\leq d_{3,\max}$, retaining its $K_3^{-2}$ dependence.
Thus \eqref{tgdc:proved_boxes} does not assert a new unrestricted
viscosity range or assume a further return.

\subsection{Signed amplitudes and first derivatives with the diffusion retained}

\begin{theorem}\label{tgdc:first_derivatives}
For every $t_c\in I_3$, both $\Theta_3(t_c)$ and $\Omega_3(t_c)$
are negative and strictly decreasing. Their exact component rates are
\begin{equation}\label{tgdc:exact_first}
 \dot\Theta_3=-P_3(a_3u-\lambda_3),\qquad
 \dot\Omega_3=-P_3(b_3-\lambda_3u).
\end{equation}
They obey the bounds
\begin{equation}\label{tgdc:component_bounds}
\begin{gathered}
 S_3e^{\gamma_i\tau/(8\sqrt2)}
 \leq P_3(t_c)\leq\min\{P_3^*,S_3e^{2\gamma_i\tau}\},\\
 \frac{u_i}{\sqrt2}P_3\leq |\Omega_3|\leq2u_iP_3,\\
 \left(\frac{\gamma_i}{4\sqrt2}-dK_3^2\overline R\right)P_3
 \leq |\dot\Theta_3|\leq2\gamma_iP_3,\\
 u_i\left(\frac{\gamma_i}{2}-2dK_3^2\overline R\right)P_3
 \leq |\dot\Omega_3|\leq\gamma_i u_iP_3.
\end{gathered}
\end{equation}
In particular both lower derivative bounds are positive, and
\begin{equation}\label{tgdc:norm_bounds}
\begin{gathered}
 P_3\sqrt{1+u_i^2/2}\leq |y|_2
       \leq P_3\sqrt{1+4u_i^2}\leq P_3^*\sqrt{1+4u_i^2},\\
 |\dot y|_2\leq\gamma_iP_3\sqrt{4+u_i^2}
       \leq\gamma_iP_3^*\sqrt{4+u_i^2}.
\end{gathered}
\end{equation}
The original beginning and attained endpoint have the exact derivatives
\begin{equation}\label{tgdc:endpoint_derivatives}
\begin{aligned}
 \dot y(t_b)&=-S_3(\gamma_i-dK_3^2)\binom1{u_i},\\
 y(t_3^a)&=-P_3^*\binom1{u(t_3^a)},\\
 \dot y(t_3^a)&=-P_3^*
 \binom{a_3(t_3^a)u(t_3^a)-dK_3^2R_3(t_3^a)}
       {b_3(t_3^a)-dK_3^2R_3(t_3^a)u(t_3^a)}.
\end{aligned}
\end{equation}
\end{theorem}
\begin{proof}
The signs and quotient in \eqref{tgdc:proved_boxes} give
$y=-P_3(1,u)^T$. Substituting this expression in the exact pair proves
\eqref{tgdc:exact_first}. Integrating the already proved logarithmic
rate from $t_b$ and using the attained endpoint and strict increase
of $P_3$ gives the first line of \eqref{tgdc:component_bounds}.
The quotient box gives its second line. Since $a_i u_i=\gamma_i$
and $b_i=\gamma_i u_i$, direct multiplication of the positive boxes
gives
\[
 \frac{\gamma_i}{4\sqrt2}\leq a_3u\leq2\gamma_i,
 \qquad
 \frac{\gamma_i u_i}{2}\leq b_3\leq\gamma_i u_i,
 \qquad 0<\lambda_3\leq dK_3^2\overline R.
\]
Consequently the displayed lower bound for $a_3u-\lambda_3$ follows.
It is at least $\gamma_i/(8\sqrt2)>0$, because
$\overline R<2$ and $dK_3^2\leq\gamma_i/(16\sqrt2)$.
Similarly,
\[
 b_3-\lambda_3u
 \geq u_i(\gamma_i/2-2dK_3^2\overline R)
 >\gamma_i u_i(1/2-1/(4\sqrt2))
 >\gamma_i u_i/4>0.
\]
Thus both derivatives in \eqref{tgdc:exact_first} are strictly negative.
Dropping their nonnegative diffusion losses only for the purpose of
an upper bound gives the two upper derivative estimates; the exact
rates continue to contain the losses. Squaring the two amplitude and
two derivative bounds, adding them, and taking the positive square
root proves \eqref{tgdc:norm_bounds}. Finally $R_3(t_b)=1$,
$u(t_b)=u_i$, $a_3(t_b)=a_i$ and $b_3(t_b)=b_i$ give the first
endpoint identity. Substitution of the proved target value gives
the other two. In particular no later endpoint was substituted for
the actual $u(t_3^a)$.
\end{proof}

\subsection{Explicit coefficient rates}

Set $\overline H=h_b+144a$, $\overline G=g_b+144b$ and
$\overline W=9\sigma_1$. Differentiating the complete coefficients
gives
\begin{equation}\label{tgdc:coefficient_rates}
\begin{aligned}
 \dot a_3&=-a_3\left(\frac{dE_3}{N_3}
                            +\frac{2bg\Omega_1}{R_3}\right),\\
 \dot b_3&=-b_3\frac{ah\Omega_1}{\chi^2},\qquad
 \dot\lambda_3=2dK_3^2bg\Omega_1.
\end{aligned}
\end{equation}
Every quantity on the right comes from the moving parent. In particular
the exact numerator derivative \eqref{tgdc:numerator} is used before
any upper estimate. Because $K_1<K_2$, $\chi^2<10$ and all summands
of $N_3$ are positive, $E_3\leq10K_2^2N_3$. Thus
\begin{equation}\label{tgdc:coefficient_rate_bounds}
\begin{gathered}
 |\dot a_3|\leq
 a_i(10dK_2^2+2b\overline G\overline W),\qquad
 |\dot b_3|\leq
 b_i a\overline H\overline W\chi_b^{-2},\\
 |\dot\lambda_3|\leq
 2dK_3^2 b\overline G\overline W.
\end{gathered}
\end{equation}
These formulas also display the distinct inherited $K_2^2$ loss and
new $K_3^2$ loss. They do not replace either by the source scaling
parameter $\nu$.

\subsection{A closed all-orders derivative calculation}

Here are finite recurrences whose values are explicit upper bounds
from the actual birth data and the already proved interval. They use
ordinary derivatives, with no change to physical time or amplitudes.
For integers $n\geq0$ and $m\geq2$ define the product expression
\[
 \mathcal P_n(X^{(1)},\ldots,X^{(m)})
 =\sum_{j_1+\cdots+j_m=n}
       \frac{n!}{j_1!\cdots j_m!}
       X^{(1)}_{j_1}\cdots X^{(m)}_{j_m}.
\]
This notation is a finite sum, and repeated arguments mean repeated
factors. Define the nonnegative arrays $H,W,A,P,Q,G,N,J$ initially by
\begin{equation}\label{tgdc:majorant_initial}
\begin{gathered}
 H_0=\overline H,\quad W_0=\overline W,\quad
 A_0^{\rm bd}=A_{1,b},\quad P_0=P_{2,b},\quad
 Q_0=\chi_b^{-1},\quad G_0=\overline G,\quad
 N_0=N_i,\quad J_0=1.
\end{gathered}
\end{equation}
In the following recurrences the array $A$ has entries $A_n^{\rm bd}$;
the superscript distinguishes its zeroth entry from the unchanged
physical constant $A_0$ in \eqref{tgdc:original}. Write
\[
 X_n=\mathcal P_n(H,H)+a^2\mathbf1_{n=0}.
\]
Having computed all eight arrays up to index $n$, compute index $n+1$
by the explicit formulas
\begin{equation}\label{tgdc:parent_majorants}
\begin{aligned}
 H_{n+1}&=aW_n,\\
 W_{n+1}&=a\mathcal P_n(A,Q)+dK_1^2W_n,\\
 A_{n+1}^{\rm bd}&=S_*W_n+dK_1^2A_n^{\rm bd},\\
 P_{n+1}&=dK_2^2\mathcal P_n(X,P),\\
 Q_{n+1}&=a\mathcal P_n(H,W,Q,Q,Q),\\
 G_{n+1}&=bW_n,\\
 N_{n+1}&=dbK_1^2A_n^{\rm bd}
                 +dcK_2^3\mathcal P_n(X,P),\\
 J_{n+1}&=2b\mathcal P_n(G,W,J,J).
\end{aligned}
\end{equation}
For the phase-length array use
\[
 R_0^{\rm bd}=\overline R,\qquad
 R_n^{\rm bd}=\mathcal P_n(G,G)\quad(n\geq1).
\]
The coefficient arrays are
\begin{equation}\label{tgdc:coefficient_majorants}
\begin{gathered}
 \mathfrak A_0=a_i,\quad \mathfrak B_0=b_i,\quad
 \mathfrak D_0=dK_3^2\overline R,\\
 \mathfrak A_n=K_3^{-1}\mathcal P_n(N,J),\qquad
 \mathfrak B_n=K_3cQ_n,\qquad
 \mathfrak D_n=dK_3^2R_n^{\rm bd}\quad(n\geq1).
\end{gathered}
\end{equation}
No optimization of these finite sums is needed for their validity.
In particular $\mathfrak A_1$ may be replaced by the smaller of its
value above and the first bound in \eqref{tgdc:coefficient_rate_bounds},
and similarly for the other two first coefficient derivatives.

The amplitude arrays start with the actual target and the sharper
first-derivative bounds already proved:
\begin{equation}\label{tgdc:amplitude_majorants}
\begin{gathered}
 T_0=P_3^*,\qquad O_0=2u_iP_3^*,\qquad
 T_1=2\gamma_iP_3^*,\qquad O_1=\gamma_i u_iP_3^*,\\
 T_{n+1}=\sum_{j=0}^n\binom nj
      (\mathfrak D_jT_{n-j}+\mathfrak A_jO_{n-j}),\qquad n\geq1,\\
 O_{n+1}=\sum_{j=0}^n\binom nj
      (\mathfrak B_jT_{n-j}+\mathfrak D_jO_{n-j}),\qquad n\geq1.
\end{gathered}
\end{equation}
The choice to begin these recurrences at $n=1$ preserves the stronger
first-derivative estimates instead of substituting a larger matrix-norm
bound.

\begin{theorem}\label{tgdc:all_orders}
For every integer $n\geq0$ the finite arrays above satisfy on $I_3$
\begin{equation}\label{tgdc:all_order_bounds}
\begin{gathered}
 |\partial_{t_c}^n\Theta_3|\leq T_n,\qquad
 |\partial_{t_c}^n\Omega_3|\leq O_n,\qquad
 |\partial_{t_c}^ny|_2\leq\sqrt{T_n^2+O_n^2},\\
 |\partial_{t_c}^na_3|\leq\mathfrak A_n,\qquad
 |\partial_{t_c}^nb_3|\leq\mathfrak B_n,\qquad
 |\partial_{t_c}^n\lambda_3|\leq\mathfrak D_n.
\end{gathered}
\end{equation}
The exact, signed amplitude derivative recurrence underlying these
bounds is
\begin{equation}\label{tgdc:exact_derivative_recurrence}
 \partial_{t_c}^{n+1}y
     =\sum_{j=0}^n\binom nj
            (\partial_{t_c}^jC_3)(\partial_{t_c}^{n-j}y),
 \qquad n\geq0.
\end{equation}
In particular the exact second derivative contains the changing
coefficient and the complete diffusion square:
\begin{equation}\label{tgdc:second_derivative}
 \ddot y=
 \begin{pmatrix}
 \lambda_3^2+a_3b_3-\dot\lambda_3&\dot a_3-2\lambda_3a_3\\
 \dot b_3-2\lambda_3b_3&\lambda_3^2+a_3b_3-\dot\lambda_3
 \end{pmatrix}y.
\end{equation}
\end{theorem}
\begin{proof}
Introduce the exact auxiliary functions $q=\chi^{-1}$ and
$j=R_3^{-1}$, both well-defined by \eqref{tgdc:proved_boxes}.
Direct differentiation of their specified inverses gives
\[
 \dot q=-ah\Omega_1q^3,\qquad
 \dot j=-2bg\Omega_1j^2.
\]
These identities retain the actual phase lengths; the symbols do not
replace them by constant values. The complete equations for
$(h,\Omega_1,A_1,P_2,q,g,N_3,j)$ are therefore the signed polynomial
system consisting of \eqref{tgdc:parent}, \eqref{tgdc:numerator} and
these two equations, with $\chi^2=h^2+a^2$ in the parent losses.
Their zero-order absolute values are bounded respectively by
$(H_0,W_0,A_0^{\rm bd},P_0,Q_0,G_0,N_0,J_0)$:
the claims for $q$ and $j$ use $\chi\geq\chi_b$ and $R_3\geq1$.

For any $m$ differentiable factors the ordinary $n$th derivative of
their product is $\mathcal P_n$ applied to their derivative arrays.
To prove this formula without assuming it, start with $n=0$.
Differentiating a term of order $n$ once chooses one factor whose
index increases. At indices summing to $n+1$, the resulting coefficient
is the sum of $n!/(j_1!\cdots(j_k-1)!\cdots j_m!)$ over indices
$k$ with $j_k>0$. It is
$n!(j_1+\cdots+j_m)/(j_1!\cdots j_m!)=(n+1)!/(j_1!\cdots j_m!)$.
This proves the formula by induction.

Apply this exact product derivative identity to the signed parent
equations and then the triangle inequality. For instance
\[
 \partial_{t_c}^{n+1}N_3
  =-dbK_1^2\partial_{t_c}^n A_1
   -dcK_2^3\partial_{t_c}^n\{(h^2+a^2)P_2\},
\]
and its absolute value is bounded by the displayed recurrence for
$N_{n+1}$. This example shows explicitly that neither inherited
diffusion term is omitted. The other seven equations give exactly
\eqref{tgdc:parent_majorants}. Induction in $n$ proves all eight parent
derivative bounds. The identities $R_3=g^2+b^2$,
$a_3=N_3j/K_3$, $b_3=K_3cq$ and $\lambda_3=dK_3^2R_3$ then prove
\eqref{tgdc:coefficient_majorants} for $n\geq1$ by the same product
formula. Their zero-order bounds were already proved in
\eqref{tgdc:proved_boxes}.

Repeatedly differentiate the exact pair to obtain
\eqref{tgdc:exact_derivative_recurrence}. Its componentwise absolute
value gives \eqref{tgdc:amplitude_majorants}, with the zero and first
orders supplied by Theorem~\ref{tgdc:first_derivatives}. Induction
now proves the two amplitude bounds at every order, and the Euclidean
norm bound follows by summing squares. Finally
$\ddot y=(\dot C_3+C_3^2)y$ and direct multiplication of the two by two
matrix give \eqref{tgdc:second_derivative}. Every recurrence uses
finitely many sums at each specified order. Thus each resulting bound
is a determinate expression in the original physical constants and
actual birth data, rather than an unevaluated supremum of a missing
coefficient derivative.
\end{proof}

\subsection{Derivative transport through the specified source clock}

Let $L_s>0$ be the source slot length used for this actual interval,
and define its specified clock by
\[
 \alpha_c=\tau_a/L_s,\qquad
 t_c=t_b+\alpha_c v,\qquad 0\leq v\leq L_s.
\]
The two actual endpoints are retained. Since $\alpha_c$ is constant
for the fixed trajectory and slot, repeated chain differentiation
gives, exactly,
\begin{equation}\label{tgdc:clock}
 \partial_v^ny(t_b+\alpha_cv)
      =\alpha_c^n\partial_{t_c}^ny(t_b+\alpha_cv),\qquad
 |\partial_v^ny|_2\leq\alpha_c^n\sqrt{T_n^2+O_n^2}.
\end{equation}
In a source chart with fixed label $Q$, exponent $h_{\rm src}$,
and $\partial_t v=Q^{-1-h_{\rm src}}$, the derivative contribution
acting on this amplitude alone is consequently
\[
 \partial_t^ny(t_b+\alpha_cv)
  =(\alpha_cQ^{-1-h_{\rm src}})^n
                      \partial_{t_c}^ny(t_b+\alpha_cv).
\]
This identity holds with the stated fixed chart parameters and affine
clock; derivatives of any additional frame, cutoff, or varying
parameter must be included by the ordinary product rule. In particular
neither the clock factor nor the physical frequency has been replaced
by one. These finite derivative estimates establish no uniform
terminal extension for an infinite sequence.

\clearpage\part{Exact clock, signed frame, and actual source pulse}
\section{The actual third-growth pair in the source clock and frame}
\label{tgframe:section}

This calculation uses equations (6.2)--(6.12), (7.2)--(7.12), and
Lemma 7.4 of the retained 166-page source. Its SHA-256 is displayed
in two consecutive segments:
\[
\begin{gathered}
\texttt{0e779481c4da40bd28d1e642e1d8ca57}\\
\texttt{447d129610df28dfa5a11e9af8ae228f}.
\end{gathered}
\]
The coupled input is the actual third-growth interval proved in
\path{finite_stage/third_growth_entry.tex}. The source's pulse
equation is written at physical viscosity one. We retain an arbitrary
positive physical viscosity $\nu_{\rm NS}$ when computing a candidate
residual; when identifying the source's particular homogeneous pulse,
we explicitly set $\nu_{\rm NS}=1$ and harmonic $m=1$.
The original coupled time scale $\nu=\sqrt{A_0}$, the physical coupled
diffusivity $d$, and the source chart parameter $\varepsilon=Q^h$
are distinct.
The source growth constant denoted $\lambda_0$ in $\lambda(v)$ below
is the source's own growth constant. It is not identified with the
coupled integer frequency $\lambda_0=\mu=\lceil\sqrt{A_0}/2\rceil$.

\subsection{Exact physical differentiation of the source pulse clock}

Fix a source label and one of its local lifts. Set
\[
 b_g=\sqrt2-1,\quad v_r=(1,-b_g),\quad v_t=(b_g,1),\quad
 J_g=\begin{pmatrix}3&1\\1&5\end{pmatrix},\quad
 T_g=4+\sqrt2,\quad\Lambda_g=4-\sqrt2.
\]
Direct multiplication gives $J_gv_t=T_gv_t$ and
$J_gv_r=\Lambda_gv_r$, while $v_t\cdot v_r=0$ and
$|v_t|^2=|v_r|^2=1+b_g^2$. The source retains
\[
\begin{gathered}
 d_r=2\left((1+h)\frac{\log\Lambda_g}{\log T_g}-h\kappa_s\right),\qquad
 \kappa_s=10^{-5},\\
 Q=2^{-\ell},\quad
 i=\left\lfloor\log_{T_g}\frac{Q^{-1-h}}{S_*}\right\rfloor,\quad
 c_i=T_g^iQ^{1+h}.
\end{gathered}
\]
On the local lift
\[
 Y_i-c_\gamma-k_{\rm copy}=\xi_gv_r+\eta_gv_t,\qquad
 v=\frac{\eta_g+r_0}{c_i},\qquad L_s=\frac{2r_0}{c_i},
 \quad
 \eta_g=\frac{v_t\cdot(Y_i-c_\gamma-k_{\rm copy})}{1+b_g^2}.
\]
The label and integer $k_{\rm copy}$ are held fixed on this lift.
Consequently $L_i\eta_g=0$, $N_i\eta_g=1$, where
$L_i=v_r\cdot\partial_{Y_i}$ and $N_i=v_t\cdot\partial_{Y_i}$.
The clock has no explicit dependence on the slow variables $R,Z,T$.
The evaluated chart derivatives are exactly
\[
 D_r=\partial_R+M_id_rR^{d_r-1}L_i,\qquad
 D_z=\varepsilon\partial_Z,\qquad
 t_*=-\varepsilon\partial_T+c_iN_i,
 \quad M_i=\Lambda_g^iQ^{d_r/2}.
\]
It follows by applying these displayed operators, without estimating
any coefficient, that
\begin{equation}\label{tgframe:clockchart}
 D_rv=D_zv=0,\qquad t_*v=1.
\end{equation}
There is also an explicit physical verification. Evaluate at the
source map $Y=v_rr^{d_r}+v_tt$ modulo the lattice. Locally the lift of
$Y_i=J_g^iY$ differs from
$\Lambda_g^iv_rr^{d_r}+T_g^iv_tt$ by a constant lattice vector.
The dual formula for $\eta_g$ therefore gives
$\eta_g=T_g^it+C_\gamma$ with a locally constant $C_\gamma$.
Thus, on the evaluated local rectangle,
\begin{equation}\label{tgframe:clockphysical}
 \partial_rv=\partial_zv=\partial_\theta v=0,\qquad
 \partial_tv=\frac{T_g^i}{c_i}=Q^{-1-h}.
\end{equation}
The $\partial_t$ in this equation is the derivative of the evaluated
physical function. On the extended domain the corresponding derivative
is $\partial_t+N_{\rm abs}$. These statements are identical after
evaluation; neither drops the auxiliary derivative.

Let $t_b$ and $r_a>0$ be the actual second return time and the actual
third target-crossing elapsed time. Define
\begin{equation}\label{tgframe:coupledclock}
 t_c(v)=t_b+\alpha v,\qquad \alpha=\frac{r_a}{L_s}>0.
\end{equation}
This parametrizes the entire proved coupled growth interval, with its
original endpoints. In physical source variables its derivative is
$\partial_tt_c=\alpha Q^{-1-h}$ and all three spatial derivatives vanish.
The positive orientation of this clock asserts no sign for a velocity.
The coupled amplitudes below are strictly negative.

\subsection{The actual frame, its inverse, and the direct candidate}

Write the original source phase normal as
\[
 n=\left(x_0-v(pF_R+p_zG_R),\ \frac pR,\
             p_z-\varepsilon v(pF_Z+p_zG_Z)\right).
\]
All $p,p_z,x_0,k,Q$ are the original fixed label constants, including
$kp\in\mathbb Z\setminus\{0\}$. Set
\[
\begin{gathered}
 K_a=\frac{n_{\tan}}{|n_{\tan}|},\quad
 N_a=((K_a)_z,-(K_a)_\theta),\quad
 s_a=\frac{n_r}{|n_{\tan}|},\\
 c(v)=c_0\sqrt{1+s(v)^2},
 \quad
 s(v)=\sigma\left(\frac{u_*}{2}+\frac{u_*v}{L_s}\right).
\end{gathered}
\]
Here $c_0<0$ is the unchanged source constant, and $K_a,N_a$
are identified with their tangential three-vectors. In particular
$c(v)\ne0$. Because $p/R\ne0$, every formula is defined throughout
the source rectangle. The actual source frame is
\[
 U=[e_r-s_aK_a,\ N_a],\qquad
 M=\begin{pmatrix}1&1\\c&-c\end{pmatrix},\qquad B=UM.
\]
The two columns of $U$ are orthogonal, with squared lengths
$g_a^2=1+s_a^2$ and $1$. Both are orthogonal to $n$.
Since $\det M=-2c\ne0$, $B$ is a bijection from $\mathbb R^2$
onto $n^\perp$. The source left inverse, defined on all of
$\mathbb R^3$, is
\begin{equation}\label{tgframe:leftinverse}
 B_\ell=M^{-1}\begin{pmatrix}e_r^T\\N_a^T\end{pmatrix},
 \qquad
 M^{-1}=\frac12\begin{pmatrix}1&c^{-1}\\1&-c^{-1}\end{pmatrix}.
\end{equation}
Indeed the last two rows applied to $U$ give $I_2$.

The eigenvalues of
\[
 B^TB=
 \begin{pmatrix}g_a^2+c^2&g_a^2-c^2\\
                 g_a^2-c^2&g_a^2+c^2\end{pmatrix}
\]
are $2g_a^2$ and $2c^2$, with orthogonal eigenvectors $(1,1)$
and $(1,-1)$. The Euclidean operator norm and the inverse norm on
the tangent plane are therefore exactly
\begin{equation}\label{tgframe:framenorms}
 \|B\|=\sqrt2\max(g_a,|c|),\qquad
 \|(B|_{\mathbb R^2})^{-1}\|_{n^\perp\to\mathbb R^2}
       =\frac1{\sqrt2\min(g_a,|c|)}.
\end{equation}
The ambient norm of the particular extension \eqref{tgframe:leftinverse}
is $(1/\sqrt2)\max(1,|c|^{-1})$; it need not equal the inverse norm
restricted to $n^\perp$.

Retain the actual coupled data
\[
 K_3=\mu\widehat\lambda_3,\quad
 S_3=\frac{A_0}{\mu}\widehat\lambda_3^{-k_3-6},\quad
 P_3^*=\frac{A_0}{\mu}\widehat\lambda_3^{-7/8}=S_3e^{L_3},
 \quad u_i=\sqrt{\frac{b_i}{a_i}}
           =\frac{K_3}{\sigma_2\sqrt{n_i}}.
\]
On the interval \eqref{tgframe:coupledclock}, the already constructed
solution is
\[
 y_3(t_c)=\binom{\Theta_3}{\Omega_3}
        =-P(v)\binom{1}{u(v)},\qquad
 S_3\le P(v)\le P_3^*,\qquad
 \frac{u_i}{\sqrt2}\le u(v)\le2u_i,
\]
with $P(0)=S_3$, $P(L_s)=P_3^*$ and $u(0)=u_i$.
These inequalities follow from the actual evolving-parent system,
rather than from a new frame approximation.
The direct candidate $t_{\rm dir}=By_3$ consequently has the exact
components
\begin{align}
 (t_{\rm dir})_r&=-P(1+u),\nonumber\\
 (t_{\rm dir})_{\tan}
       &=P s_a(1+u)K_a-Pc(1-u)N_a,\label{tgframe:directcomponents}\\
 |t_{\rm dir}|^2
       &=P^2\bigl[(1+s_a^2)(1+u)^2+c^2(1-u)^2\bigr],\nonumber\\
 (t_{\rm dir})_r(t_{\rm dir})_{\tan}
       &=-P^2s_a(1+u)^2K_a+P^2c(1-u^2)N_a.\nonumber
\end{align}
These follow by multiplying $My_3=-P(1+u,c(1-u))^T$ and using
the orthogonality of the two columns of $U$. In particular
\[
 |(t_{\rm dir})_r|\ge S_3(1+u_i/\sqrt2),\qquad
 |t_{\rm dir}|\le\sqrt2\max(g_a,|c|)P_3^*\sqrt{1+4u_i^2}.
\]
At the actual target,
\begin{equation}\label{tgframe:endpointomega}
 \frac{P_3^*u_i}{\sqrt2}\le|\Omega_3(t_b+r_a)|
                      \le2P_3^*u_i,\qquad
 P_3^*u_i=\frac{A_0}{\sigma_2\sqrt{n_i}}
                          \widehat\lambda_3^{1/8}.
\end{equation}
The last equality retains $\mu$ until it cancels against its identical
factor in $K_3$. Thus a decreasing temperature target does not give a
decreasing two-component target.

\subsection{The complete principal mismatch, including the cutoff}

With a prime denoting $v$ differentiation at fixed slow variables,
retain
\[
 \mathsf K=\begin{pmatrix}0&-2F&0\\2F+RF_R&0&0\\G_R&0&0\end{pmatrix},
 \quad
 A_\Phi=-\mathsf K+
        \frac{n(n^T\mathsf K-(n')^T)}{|n|^2},
 \quad H=B_\ell(A_\Phi B-B').
\]
Differentiating $n^TB=0$ shows
$n^TB'=-(n')^TB$. Direct multiplication gives
$n^TA_\Phi B=-(n')^TB$ as well. Therefore
$A_\Phi B-B'$ takes values in $n^\perp$, and
\begin{equation}\label{tgframe:frametransport}
 A_\Phi B-B'=BH.
\end{equation}
The source defines the exact error matrix
$E=H-\operatorname{diag}(\lambda,-\lambda)$ with
$\lambda=\lambda_0/\sqrt{1+s^2}$. Its source estimate
$\|E\|\le C/S_*$ is not an assertion that $E$ is zero.

The actual coupled matrix and arbitrary-viscosity source loss are
\[
\begin{gathered}
 C_c=\begin{pmatrix}-d_c&a_3\\b_3&-d_c\end{pmatrix},
 \quad d_c=dK_3^2R_3,\quad
 a_3=\frac{N_3}{K_3R_3},\quad b_3=\frac{K_3c_{\rm det}}{\chi_{\rm old}},\\
 \quad
 \delta_\nu=\nu_{\rm NS}m^2\varepsilon k^2|n|^2,\quad
 \Delta_\nu=\delta_\nu-\alpha d_c.
\end{gathered}
\]
Here $c_{\rm det}$ and $\chi_{\rm old}$ are the third-growth determinant
constant and older phase length; these names distinguish them from the
source $c(v)$ and the cutoff below. Every coupled quantity is evaluated
at $t_c(v)$. The actual third growth has zero added amplitude controls,
so $y_3'=\alpha C_cy_3$ exactly.

Let $\chi=\eta(R,Z,T)\chi_g(\xi_g)\psi(v)$ be a smooth scalar
cutoff with the source support, and set $t_\chi=\chi By_3$.
The homogeneous-pressure formula, with zero principal source, is
\[
 \pi_\chi=\frac{i}{km|n|^2}
        \bigl(n\cdot\mathsf Kt_\chi-n'\cdot t_\chi\bigr).
\]
The cutoff leaves $n^Tt_\chi=0$. Differentiating this identity
and substituting the pressure gives
\begin{equation}\label{tgframe:residual}
 t_\chi'+\mathsf Kt_\chi+\delta_\nu t_\chi+ikmn\pi_\chi
 =B\mathcal E_\chi,\qquad
 \mathcal E_\chi=\chi(\alpha C_c-H+\delta_\nu I)y_3+\chi'y_3.
\end{equation}
For completeness, the normal component of the sum before pressure is
$-n'\cdot t_\chi+n\cdot\mathsf Kt_\chi$; multiplication by $ikmn$
times the displayed pressure cancels it. The remaining vector is
$t_\chi'-A_\Phi t_\chi+\delta_\nu t_\chi$. Use
\eqref{tgframe:frametransport} and $y_3'=\alpha C_cy_3$ to obtain
\eqref{tgframe:residual}.

In original scalar entries this is
\begin{equation}\label{tgframe:residualentries}
 \mathcal E_\chi=-P
 \left\{\chi
 \binom{\Delta_\nu-\lambda-E_{11}+(\alpha a_3-E_{12})u}
       {\alpha b_3-E_{21}+(\Delta_\nu+\lambda-E_{22})u}
       +\chi'\binom1u\right\}.
\end{equation}
All entries are signed. In particular the damping mismatch is
$\delta_\nu-\alpha dK_3^2R_3$, not the sum of the losses.
At fixed slow variables and fixed $\xi_g$, $\chi'=\eta\chi_g\psi'$.
The physical derivative of $\chi$ also includes
$-\varepsilon(\partial_T\eta)\chi_g\psi$ in $t_*\chi$; that slow term
belongs to the complete physical residual and has not been included in,
or removed by, the principal equation \eqref{tgframe:residual}.
A direct finite bound is
\[
 |\mathcal E_\chi|
 \le P\sqrt{1+u^2}\left[
 |\chi|\bigl(|\Delta_\nu|+\lambda+\|E\|
                  +\alpha\max(a_3,b_3)\bigr)+|\chi'|\right].
\]
This is the triangle inequality and the exact norm
$\left\|\begin{smallmatrix}0&a_3\\b_3&0\end{smallmatrix}\right\|
=\max(a_3,b_3)$, with $a_3,b_3>0$.
It retains the carrier, frame, cutoff and all physical diffusion factors.

\subsection{An explicit map of the actual seed to the actual source pulse}

For this subsection only, take $m=1$, $\nu_{\rm NS}=1$, so
$\delta_1=\varepsilon k^2|n|^2$. This is precisely the source
homogeneous pulse equation with its given background. Put
\[
 d_{\rm ref}=\varepsilon k^2B_s^2(1+s^2),\qquad
 {\cal P}_\sigma(v)=
   \exp\left(\int_{L_s/2}^{v}(\lambda(w)-d_{\rm ref}(w))\,dw\right).
\]
Every coefficient is finite on the closed interval. Thus
$\mathcal P_\sigma(0)>0$ exactly, without a Gaussian approximation.
Source Lemma 7.4 specifies the initial coordinate vector
\[
 z_{\sigma,0}=\binom{{\cal P}_\sigma(0)}0,\qquad
 t^h_\sigma(0)=B(0)z_{\sigma,0}.
\]
The fixed source label may affect this strictly positive scalar.

For any nonzero real target vector $z_0=(z_1,z_2)^T$, define
\begin{equation}\label{tgframe:L0general}
 L_0=-\frac1{S_3(1+u_i^2)}
 \begin{pmatrix}z_1&-z_2\\z_2&z_1\end{pmatrix}
 \begin{pmatrix}1&u_i\\-u_i&1\end{pmatrix}.
\end{equation}
This is an explicit invertible linear map with
\[
 L_0[-S_3(1,u_i)^T]=z_0,\qquad
 \det L_0=\frac{|z_0|^2}{S_3^2(1+u_i^2)}>0,
\]
and inverse
\begin{equation}\label{tgframe:L0inverse}
 L_0^{-1}=-\frac{S_3}{|z_0|^2}
       \begin{pmatrix}1&-u_i\\u_i&1\end{pmatrix}
       \begin{pmatrix}z_1&z_2\\-z_2&z_1\end{pmatrix}.
\end{equation}
Indeed the second matrix of \eqref{tgframe:L0general} sends
$(1,u_i)^T$ to $(1+u_i^2,0)^T$, and each matrix times its
transpose is a positive scalar times $I_2$.
The same multiplication verifies the inverse and proves the exact norms
\begin{equation}\label{tgframe:L0norms}
 \|L_0\|=\frac{|z_0|}{S_3\sqrt{1+u_i^2}},\qquad
 \|L_0^{-1}\|=\frac{S_3\sqrt{1+u_i^2}}{|z_0|}.
\end{equation}
In particular the source target gives
\[
 L_0=-\frac{{\cal P}_\sigma(0)}{S_3(1+u_i^2)}
        \begin{pmatrix}1&u_i\\-u_i&1\end{pmatrix},
 \quad
 \|L_0\|=
 \frac{{\cal P}_\sigma(0)\,\mu\widehat\lambda_3^{k_3+6}}
 {A_0\sqrt{1+K_3^2/(\sigma_2^2n_i)}}.
\]
No seed or source initial amplitude has been set to one.

Define the two identity-initial-value fundamental matrices by
\[
 X_c'=\alpha C_cX_c,\quad X_c(0)=I_2,\qquad
 X_s'=(H-\delta_1I_2)X_s,\quad X_s(0)=I_2.
\]
Here is a direct existence, invertibility and uniqueness construction.
For a continuous matrix $D$ on $[0,L_s]$ let
\[
 X_D(v)=I_2+
 \sum_{j\ge1}\int_{0\le w_j\le\cdots\le w_1\le v}
       D(w_1)\cdots D(w_j)\,dw_j\cdots dw_1 .
\]
If $\sup\|D\|=M_D$, the $j$th term is bounded by
$(M_DL_s)^j/j!$. Uniform convergence and regrouping the integrals
give $X_D=I_2+\int_0^vD(w)X_D(w)\,dw$, hence the differential
equation. The corresponding right-ordered series for
$Z'=-ZD$, $Z(0)=I_2$, converges by the same bound.
Differentiation gives $(ZX_D)'=0$, hence $ZX_D=I_2$.
For square matrices this is invertibility. If two solutions agree
initially, the difference satisfies the zero-initial-value integral
equation; iteration bounds its norm by a fixed finite bound times
$(M_Dv)^j/j!$ for every $j$, which tends to zero.
This proves uniqueness. Smooth coefficient dependence follows from
the same series with the following explicit bound. On a compact
coefficient-parameter neighborhood let $M_N\ge1$ bound every coefficient
derivative through order $N$. Differentiating a product of $j$ factors
$N$ times allocates each derivative to one of $j$ factors, so there are
at most $j^N$ terms, each of norm at most $M_N^j$. The differentiated
$j$th integral is therefore bounded by
$j^N(M_NL_s)^j/j!$, a summable sequence. Uniform convergence of these
derivative series justifies termwise differentiation: integrate each
first-derivative series along a parameter segment and interchange the
uniform sum with that integral, then iterate through order $N$.
The right-ordered inverse series has the same bound. Pulse derivatives
then follow from the differential equations. This proves smoothness,
with the explicit derivative equations below providing further bounds.

Now set
\begin{equation}\label{tgframe:intertwiner}
 L(v)=X_s(v)L_0X_c(v)^{-1},\qquad
 L^{-1}=X_cL_0^{-1}X_s^{-1}.
\end{equation}
Differentiation gives
\[
 L'=(H-\delta_1I_2)L-\alpha LC_c.
\]
The original solution is $y_3(v)=X_c(v)y_3(0)$.
Consequently $L(v)y_3(v)=X_s(v)z_{\sigma,0}$, and
\begin{equation}\label{tgframe:exactpulsematch}
 B(v)L(v)y_3(v)=t^h_\sigma(v)
\end{equation}
by uniqueness of the source homogeneous equation and the exact initial
seed calculation. Thus this finite map reaches the actual source pulse.
It does not merely send the coupled pair into an unspecified transverse
solution. Multiplying both sides by the identical source cutoffs
preserves equality of the amplitudes, and applying identical linear
curl and pressure constructions preserves equality of the resulting
fields. Multiplication by the source prescribed scalar leading-wave
amplitude also preserves that equality. The cutoff principal residual
is then $\chi'BLy_3$; the remaining source construction uses its own
estimates for that exact source object.

\subsection{All scale costs of the finite intertwiner}

Here and below $\|\cdot\|$ is the Euclidean matrix norm.
Put $D_i=\operatorname{diag}(1,u_i)$,
$\kappa_i=\max(u_i,u_i^{-1})$, and
\[
 \beta_c(v)=\frac12\left(a_3(t_c(v))u_i+
                                  \frac{b_3(t_c(v))}{u_i}\right).
\]
The actual coefficient bounds $a_3\le a_i$, $b_3\le b_i$ and
$a_iu_i=b_i/u_i=\gamma_i$ give $0<\beta_c\le\gamma_i$.
Conjugating the coupled matrix by the constant $D_i$ gives
\[
 D_i^{-1}C_cD_i=
 \begin{pmatrix}-d_c&a_3u_i\\b_3/u_i&-d_c\end{pmatrix}.
\]
Its symmetric part has eigenvalues $-d_c\pm\beta_c$.
For any solution $w$ in these coordinates, differentiating $|w|^2$
therefore gives
\[
 2\alpha(-d_c-\beta_c)|w|^2
 \le (|w|^2)'
 \le2\alpha(-d_c+\beta_c)|w|^2.
\]
Multiplication by the corresponding scalar integrating factors proves
the following bounds, including the lower singular-value estimate
needed for the inverse:
\[
 \|X_c(v)\|\le\kappa_i
       e^{\int_0^v\alpha(\beta_c-d_c)\,dw},\qquad
 \|X_c(v)^{-1}\|\le\kappa_i
       e^{\int_0^v\alpha(\beta_c+d_c)\,dw}.
\]
For $D_s=H-\delta_1I_2=
\operatorname{diag}(\lambda,-\lambda)-\delta_1I_2+E$,
its symmetric part is bounded between
$(-\lambda-\delta_1-\|E\|)I_2$ and
$(\lambda-\delta_1+\|E\|)I_2$. The identical energy argument gives
\[
 \|X_s(v)\|\le
 e^{\int_0^v(\lambda-\delta_1+\|E\|)\,dw},\qquad
 \|X_s(v)^{-1}\|\le
 e^{\int_0^v(\lambda+\delta_1+\|E\|)\,dw}.
\]
Combining these four proved estimates with \eqref{tgframe:intertwiner}
yields
\begin{align}
 \|L(v)\|&\le\kappa_i\|L_0\|
 e^{\int_0^v[\lambda-\delta_1+\|E\|
                       +\alpha(\beta_c+d_c)]\,dw},\nonumber\\
 \|L(v)^{-1}\|&\le\kappa_i\|L_0^{-1}\|
 e^{\int_0^v[\lambda+\delta_1+\|E\|
                       +\alpha(\beta_c-d_c)]\,dw}.
 \label{tgframe:intertwinernorms}
\end{align}
All quantities on the right are the retained, finite coefficient paths.
In particular the exact coupled diffusion integral at the endpoint is
\[
 \int_0^{L_s}\alpha d_c\,dv
 =dK_3^2\int_0^{r_a}R_3(t_b+r)\,dr
 \in[dK_3^2r_a,\ 2dK_3^2r_a].
\]
The equality is the substitution $r=\alpha v$; the interval is the
proved actual phase-length bound. The determinant is also exact:
\begin{equation}\label{tgframe:detL}
 \det L(v)=\frac{|z_{\sigma,0}|^2}{S_3^2(1+u_i^2)}
 \exp\left(\int_0^v[
        \operatorname{tr}E-2\delta_1+2\alpha d_c]\,dw\right).
\end{equation}
For a differentiable $2$ by $2$ matrix, expanding its two-column
determinant and inserting $X'=DX$ gives
$(\det X)'=(\operatorname{tr}D)\det X$. Its initial determinant
is one. Apply this calculation to $X_s$ and $X_c$, and use
$\operatorname{tr}H=\operatorname{tr}E$ and
$\operatorname{tr}C_c=-2d_c$ to prove \eqref{tgframe:detL}.

Higher pulse derivatives have the explicit recurrence
\[
 L^{(n+1)}=
 \sum_{j=0}^n\binom nj\left[
 D_s^{(j)}L^{(n-j)}-L^{(n-j)}D_c^{(j)}\right],
 \qquad D_c=\alpha C_c .
\]
For a coefficient parameter derivative $\partial_\xi$ the exact
inhomogeneous equation is
\[
 (\partial_\xi L)'=D_s\partial_\xi L
       -(\partial_\xi L)D_c
       +(\partial_\xi D_s)L-L(\partial_\xi D_c).
\]
Its initial value is the derivative of the explicit
\eqref{tgframe:L0general}. To see directly that this produces a finite
bound, multiply on the left by $X_s^{-1}$ and on the right by $X_c$,
differentiate, and integrate the last two terms with these same
multipliers. The preceding norm estimates bound every resulting
integral on $[0,L_s]$. Iteration proves the analogous assertion for
every finite mixed derivative; all product-rule terms and initial
derivatives are retained. This also proves smoothness on each compact
source slow neighborhood where the frame is defined.

The factor $\mathcal P_\sigma(0)/S_3$ and the two explicit propagation
costs in \eqref{tgframe:intertwinernorms} can be very large or very small.
There is no supplied relation identifying $L_s$ with $L_3$.
Thus these equations prove a finite invertible map and exact matching
to the source object, while exposing the cost of that map. They do not
prove a uniform weighted-class estimate for $L$ or for the direct
candidate $L=I$. The source's particular pulse retains its own source
estimates because of the equality \eqref{tgframe:exactpulsematch}.
Transport of the full viscosity-one witness to another physical
viscosity uses the already proved affine orthogonal and viscosity map.
Alternatively one may replace $\delta_1$ by $\delta_\nu$ in this
finite ODE construction, but the resulting changed-background
homogeneous pulse is not asserted to be the source's published witness.

\clearpage\part{Complete physical curl, moments, and force costs}
\section{Actual third-growth budgets in the physical curl and force}
\label{tgq:section}
We now use the actual third-growth pair and the original source frame
of the preceding calculations, with the direct coefficient choice
$L=I_2$. This fixes a particular candidate within the earlier general
bridge. All source band, box, sign and nonzero harmonic indices below
are retained, and only a finite collection is used on the present compact
interior patch. The cutoff of each term has compact support inside its
source rectangle and active annulus, with smooth extension by zero.
The original coupled datum, its selected return time $t_b$ and its
growth duration $r_a$ are fixed before assigning these source labels;
they have no dependence on the source spatial or auxiliary variables.
The source time and the coupled time have the exact dictionary proved
above:
\[
 t_{c,j}=t_b+\alpha_jv_j,\qquad
 \alpha_j=\frac{r_a}{L_{s,j}},\qquad
 \nabla_x t_{c,j}=0,\qquad
 \partial_t t_{c,j}=\alpha_jQ_j^{-1-h}.
\]
These derivative identities concern the complete physical phase
evaluation on each lifted rectangle. They remain valid when differentiating
a field supported strictly inside that rectangle: the field vanishes on
a collar where its extension would otherwise be needed.

Write $P(t_c)=-\Theta_3(t_c)$, $u(t_c)=\Omega_3(t_c)/\Theta_3(t_c)$.
The actual finite third growth gives
\[
 S_3\le P\le P_3^*,\qquad
 \frac{u_i}{\sqrt2}\le u\le2u_i,\qquad
 u_i=\sqrt{b_i/a_i}=\frac{K_3}{\sigma_2\sqrt{n_i}},
 \qquad \gamma_i=\sqrt{a_ib_i}.
\]
The original quantities remain
\[
 K_3=\mu\widehat\lambda_3,\qquad
 S_3=\frac{A_0}{\mu}\widehat\lambda_3^{-k_3-6},\qquad
 P_3^*=\frac{A_0}{\mu}\widehat\lambda_3^{-7/8},\qquad
 \nu=\sqrt{A_0}.
\]
In particular the following two component budgets have different
original frequency powers:
\begin{equation}
 \begin{split}
 Y_1&=P_3^*=\frac{A_0}{\mu}\widehat\lambda_3^{-7/8},\\
 Y_2&=2u_iP_3^*
      =\frac{2A_0}{\sigma_2\sqrt{n_i}}\widehat\lambda_3^{1/8},\\
 V_1&=2\gamma_iP_3^*,\qquad
 V_2=\gamma_i u_iP_3^*
      =\frac{\gamma_i A_0}{\sigma_2\sqrt{n_i}}
                       \widehat\lambda_3^{1/8}.
 \end{split}\label{tgq:budgets}
\end{equation}
Here the indices $1,2$ mean the two amplitude components, not physical
vector directions or radial moment weights. The derivative calculation
above proves
$|y_a(t_c)|\le Y_a$ and $|\dot y_a(t_c)|\le V_a$ on the entire
original interval, including its original negative seed.
For the second component at the actual threshold, the sharper two-sided
conclusion is
\[
 \frac{A_0}{\sigma_2\sqrt{2n_i}}\widehat\lambda_3^{1/8}
 \le |\Omega_3(t_3^a)|
 \le \frac{2A_0}{\sigma_2\sqrt{n_i}}\widehat\lambda_3^{1/8}.
\]
All these identities follow by multiplying the original quotient bounds
by $P$ and using $K_3=\mu\widehat\lambda_3$; the factor $\mu$ is
canceled only in the explicitly displayed product $u_iP_3^*$, not
removed from the underlying objects or their definitions.

\subsection{Two complete spatial shape fields for each retained harmonic}
\label{tgq:shapes}
For harmonic $j$, retain $k_jm_j\ne0$, where $k_j$ is the source
carrier and $m_j$ its integer harmonic. Neither equals the coupled
integer $k_3$ by definition. Let $B_j$ be the source frame, $n_j$ its
normal, and $\chi_j$ the full real slow, transverse and pulse cutoff.
All are the original specified functions on that label. Put, for
$a=1,2$ and the standard vectors $e_a\in\mathbb R^2$,
\begin{equation}
 \begin{split}
 T_{ja}&=\chi_jB_je_a,\qquad
 C_{ja}=\frac{i\,n_j\times T_{ja}}{k_jm_j|n_j|^2},\\
 W_{ja}&=\operatorname{Re}
       \bigl[(T_{ja}+\operatorname{curl}_{*,j,\mathrm{coeff}}C_{ja})
                              e^{ik_jm_j\Phi_j}\bigr],\\
 \varpi_{ja}
  &=\operatorname{Re}\left[
    \frac{i\{n_j^T\mathsf K_jT_{ja}-(n'_j)^TT_{ja}\}}
                {k_jm_j|n_j|^2}\,e^{ik_jm_j\Phi_j}\right],\\
 Z_{ja}&=Q_j^{-A}W_{ja},\qquad
 \Pi_{ja}=Q_j^{-2A}\varpi_{ja},\qquad A=\tfrac12+h .
 \end{split}\label{tgq:shape_def}
\end{equation}
Vectors in a cross product have their cylindrical components in the
oriented orthonormal physical basis. The coefficient curl is the full
operator from the preceding moment calculation. Equivalently,
\[
 Z_{ja}=\operatorname{Re}\operatorname{curl}_x
      \bigl[Q_j^{1/2-A}C_{ja}e^{ik_jm_j\Phi_j}\bigr].
\]
The transverse triple product proves this equivalence, including every
cutoff derivative. Thus each $Z_{ja}$ is a smooth compactly supported
divergence-free physical vector field; the same divergence identity
holds before auxiliary evaluation, with the commuting lifted
derivatives. The pressure coefficient is linear in $T_{ja}$ and includes
the full moving-normal term.

The actual added physical fields, with the two original amplitudes,
are exactly
\begin{equation}
 w_y=\sum_j\sum_{a=1}^2 y_a(t_{c,j}) Z_{ja},\qquad
 \pi_y=\sum_j\sum_{a=1}^2 y_a(t_{c,j})\Pi_{ja}.
 \label{tgq:factorization}
\end{equation}
Indeed $\nabla_xt_{c,j}=0$ gives
$\operatorname{curl}(y_a(t_{c,j})\mathcal A)=
y_a(t_{c,j})\operatorname{curl}\mathcal A$. It also gives
$\nabla(y_a(t_{c,j})\Pi)=y_a(t_{c,j})\nabla\Pi$.
The pressure formula in \eqref{tgq:shape_def} is linear, so it agrees
with the selected pressure for the amplitude $\chi_jB_jy(t_{c,j})$.
This proves \eqref{tgq:factorization}; in particular it does not discard
a spatial derivative of a rapidly growing amplitude. That derivative
has been calculated and is zero in these precise source coordinates.

The exact spatial derivative is
\[
 \partial_zw_y=\sum_{j,a}Q_j^{-A-D}y_a(t_{c,j})
          \partial_{Z_j}W_{ja},\qquad D=\tfrac12-h,
\]
where $\partial_{Z_j}W_{ja}$ differentiates its coefficient,
its exponential phase and all source frame and cutoff factors.
The absolute auxiliary evaluation map is independent of $z$.
There is no derivative of the actual coupled time in this formula.

\subsection{Quantitative radial moments with the original component powers}
\label{tgq:moment_budgets}
For a physical field $f$, use the precise norm
\[
 \|f\|_{[e]}^2=\int_0^\infty r^e\langle |f|^2\rangle\,dr,
 \qquad e=1,2,
\]
with the normalized angular and auxiliary averages of the moment
calculation. In a fixed band let
$\|F\|_{[e],j}^2=\int R_j^e\langle|F|^2\rangle\,dR_j$.
All norms are at the corresponding fixed slow axial coordinate and
physical time. Define the actual finite shape budgets
\begin{equation}
 \begin{split}
 \mathcal V_e&=
  \sum_{j,a}Y_a Q_j^{(e+1)/4-A}\|W_{ja}\|_{[e],j},\\
 \mathcal Z_e&=
  \sum_{j,a}Y_a Q_j^{(e+1)/4-A-D}
                              \|\partial_{Z_j}W_{ja}\|_{[e],j}.
 \end{split}\label{tgq:shape_budgets}
\end{equation}
The factors $Y_1,Y_2$ are the two explicit original-frequency expressions
in \eqref{tgq:budgets}, not a common unspecified constant. Triangle
inequality, the pointwise amplitude bounds, and
$r^e\,dr=Q_j^{(e+1)/2}R_j^e\,dR_j$ prove
\[
 \|w_y\|_{[e]}\le\mathcal V_e,\qquad
 \|\partial_zw_y\|_{[e]}\le\mathcal Z_e .
\]
No orthogonality of distinct harmonics has been assumed in these bounds.
One may evaluate every norm by the explicit source fields
\eqref{tgq:shape_def}; compactness makes each finite.

Let $w_*=u_*-\langle u_*\rangle_\theta$ be the actual old wave,
including its curl corrections, and set
\[
 \mathcal W_e=\|w_*\|_{[e]},\qquad
 \mathcal W_{e,z}=\|\partial_zw_*\|_{[e]},
\]
restricting the old-field radial integrals to a fixed enclosing support
of the new fields and their derivatives. The exact moments already
proved are
\begin{align*}
 M_2&=\partial_z\int r^2
   \langle w_{*,z}w_{y,\theta}+w_{y,z}w_{*,\theta}
                         +w_{y,z}w_{y,\theta}\rangle\,dr,\\
 M_1&=\partial_z\int r
   \langle2w_{*,z}w_{y,z}+w_{y,z}^2\rangle\,dr .
\end{align*}
Apply the product rule to every term and then Cauchy--Schwarz on its
original weighted product measure. This gives the concrete simultaneous
bounds
\begin{equation}
 |M_e|\le
 2\bigl(\mathcal W_{e,z}\mathcal V_e+
         \mathcal W_e\mathcal Z_e+
         \mathcal V_e\mathcal Z_e\bigr),\qquad e=1,2.
 \label{tgq:moment_bound}
\end{equation}
For $M_2$, the first two cross products contribute one term each to
each of the first two summands, and the self product contributes two
terms to the last. For $M_1$ the explicit coefficient $2$ and the
derivative of the square give exactly those same three coefficients.
Thus \eqref{tgq:moment_bound} retains the complete cross-and-quadratic
dependence on the original component budgets.

\subsection{The full physical residual with all clock and diffusion costs}
\label{tgq:full_residual}
Fix the actual finite source state $(u_*,p_*)$ used in the bridge and
an arbitrary physical viscosity $\nu_{\rm NS}>0$ in the residual.
It remains distinct from the coupled $d$ and $\nu=\sqrt{A_0}$.
Define the explicitly differentiated shape residual
\[
 \mathcal K_{ja}=
 \partial_tZ_{ja}+(u_*\cdot\nabla)Z_{ja}
 +(Z_{ja}\cdot\nabla)u_*-\nu_{\rm NS}\Delta Z_{ja}+\nabla\Pi_{ja}.
\]
With $C_3(t_c)=
\begin{pmatrix}-dK_3^2R_3&a_3\\ b_3&-dK_3^2R_3\end{pmatrix}$,
the full residual increment of \eqref{tgq:factorization} is
\begin{equation}
 \begin{split}
 \Delta R_y={}&
 \sum_{j,a} y_a(t_{c,j})\mathcal K_{ja}\\
 &+\sum_{j,a}\alpha_jQ_j^{-1-h}
                   (C_3y)_a(t_{c,j})Z_{ja}\\
 &+\sum_{j,a,k,b} y_a(t_{c,j})y_b(t_{c,k})
                                  (Z_{ja}\cdot\nabla)Z_{kb}.
 \end{split}\label{tgq:physical_residual}
\end{equation}
To prove it, substitute \eqref{tgq:factorization} into the complete
Navier--Stokes residual and subtract $R_{\nu_{\rm NS}}(u_*,p_*)$.
The time product rule gives
$\partial_ty_a=\alpha_jQ_j^{-1-h}(C_3y)_a$.
Every spatial derivative of $y_a(t_{c,j})$ vanishes, so neither
advection nor diffusion nor pressure adds an unrecorded derivative
of that factor. The three linear spatial terms, the time derivative
of $Z_{ja}$ and the pressure gradient are precisely $\mathcal K_{ja}$.
The remaining quadratic product is the last full sum in
\eqref{tgq:physical_residual}. This proves the identity with every
old-wave interaction, cutoff, curl term and pressure included.
Distinct labels with disjoint source supports have zero corresponding
products; no such cancellation is used for harmonics of the same label.
The baseline $R_{\nu_{\rm NS}}(u_*,p_*)$ must be added back to recover
the candidate's total residual.

The source units make every factor in this identity explicit.
Let $\nabla_{*,j}$ and $\Delta_{*,j}$ be the full cylindrical spatial
operators, including basis derivatives and the lifted auxiliary
derivatives, with
$\nabla_x=Q_j^{-1/2}\nabla_{*,j}$ and
$\Delta_x=Q_j^{-1}\Delta_{*,j}$ on that chart.
Put $U_{*,j}=Q_j^Au_*$ and
$\mathfrak t_{*,j}=Q_j^{1+h}\partial_t$ with the full physical
chain rule. Its extended expression is
$-\varepsilon_j\partial_{T_j}+c_{i,j}N_{i,j}$ as in source (6.6).
Define
\[
 \mathfrak K_{ja}=
 \mathfrak t_{*,j}W_{ja}
 +(U_{*,j}\cdot\nabla_{*,j})W_{ja}
 +(W_{ja}\cdot\nabla_{*,j})U_{*,j}
 -\nu_{\rm NS}\varepsilon_j\Delta_{*,j}W_{ja}
 +\nabla_{*,j}\varpi_{ja}.
\]
Each term is obtained by differentiating the explicit shapes in
\eqref{tgq:shape_def}, and
\[
 \mathcal K_{ja}=Q_j^{-2A-1/2}\mathfrak K_{ja}.
\]
For the time term the exponent identity is $A+1+h=2A+1/2$;
for viscosity it is
$Q_j^{-A-1}=Q_j^{-2A-1/2}\varepsilon_j$.
The nonlinear and pressure terms have the same stated residual unit.
This proves every factor rather than replacing the source or coupled
viscosity by a common normalized number.

Using Euclidean vector norms and the induced norm of the differentiated
vector map, \eqref{tgq:physical_residual} and
\eqref{tgq:budgets} give the following pointwise bound in actual
shape fields:
\begin{equation}
 \begin{split}
 |\Delta R_y|\le{}&
 \sum_{j,a}Q_j^{-2A-1/2}
       \bigl(Y_a|\mathfrak K_{ja}|+\alpha_jV_a|W_{ja}|\bigr)\\
 &+\sum_{j,a,k,b}Y_aY_bQ_j^{-A}Q_k^{-A-1/2}
                   |W_{ja}|\,\|\nabla_{*,k}W_{kb}\|.
 \end{split}\label{tgq:residual_bound}
\end{equation}
Both $Q_j$ and $Q_k$ remain in every mixed-band term.
The expression $\nabla_{*,k}W_{kb}$ includes the cylindrical frame
connection, so the bound controls the actual Cartesian directional
derivative. The clock satisfies the explicit original-stage bounds
\[
 \frac{L_3}{2\gamma_iL_{s,j}}\le\alpha_j
       \le\frac{8\sqrt2L_3}{\gamma_iL_{s,j}}.
\]
Consequently its derivative costs in the first line can, if desired,
be bounded using the original constants by
\[
 \alpha_jV_1\le
  \frac{16\sqrt2L_3}{L_{s,j}}
       \frac{A_0}{\mu}\widehat\lambda_3^{-7/8},\qquad
 \alpha_jV_2\le
  \frac{8\sqrt2L_3}{L_{s,j}}
       \frac{A_0}{\sigma_2\sqrt{n_i}}\widehat\lambda_3^{1/8}.
\]
These are finite quantitative costs for the actual inserted third stage.
They do not assert decay of the resulting residual at a singular limit.

The source pulse equation (7.5) itself belongs to its viscosity-one
source object. At a general $\nu_{\rm NS}$ in the displayed candidate
residual, the principal wave damping is
$\nu_{\rm NS}m_j^2\varepsilon_jk_j^2|n_j|^2$.
Alternatively, the exact viscosity transport already proved can be
applied to the complete viscosity-one source field and this complete
candidate. That transport rescales support, force, residual and
vorticity together. Merely changing this one damping coefficient
does not assert that the unchanged source background is a solution
with the same force at the changed viscosity.

Equations \eqref{tgq:shape_def}, \eqref{tgq:factorization},
\eqref{tgq:moment_bound} and \eqref{tgq:physical_residual} are the
explicit substitution of the actual third-growth amplitudes into
the source-coordinate construction. Every surviving moment and
residual has an actual differentiated field and a retained-scale
bound. The exact intertwiner above supplies the relation to the
source's homogeneous pulse; neither this direct candidate nor that
finite coordinate identity is silently promoted to a completed
infinite modified coupled construction.

\end{document}
